% Media and Communication — Anglais 1re Tech — Cours signé OnBuch+
% Programme officiel MINESEC. Compiler avec : tectonic N05-media-and-communication.tex
\newcommand{\DOCMATIERE}{Anglais}
\newcommand{\DOCNIVEAU}{1re Tech}
\newcommand{\DOCMODULE}{Anglais – classe de Première technique}
\newcommand{\DOCLECON}{N05}
\newcommand{\DOCTITRE}{Media and Communication}
% OnBuch+ — préambule « cours signés » ENRICHI (compilé avec tectonic / XeLaTeX)
% Système visuel « Pop » d'OnBuch+ : fond crème (jamais blanc), bordures encre
% épaisses, ombres « dures » décalées, titres Archivo Black en capitales.
% Version MATHÉMATIQUES : graphes (pgfplots/tikz), boîtes pédagogiques
% (définition, propriété, méthode, attention, le savais-tu, exercice résolu,
% exercice, TP, bilan), page de garde et signature OnBuch+ ancrée en bas.
%
% Macros attendues dans main.tex AVANT \input{preamble} :
%   \DOCMATIERE  \DOCNIVEAU  \DOCMODULE  \DOCLECON (numéro)  \DOCTITRE
\documentclass[11pt]{article}

\usepackage{fontspec}
\usepackage[french]{babel}
\usepackage[a4paper,margin=2cm,top=3.4cm,bottom=2.8cm,headheight=1.4cm,footskip=1.3cm]{geometry}
\usepackage{amsmath,amssymb,amsfonts}
\usepackage{mathtools} % \xmapsto, \coloneqq... souvent utilisés par les modèles
\usepackage{cancel} % simplifications barrées (bilans de forces)
% \abs{z} (module, valeur absolue) et \norm{u} (norme), utilisés par les modèles
\providecommand{\abs}{}\let\abs\relax\DeclarePairedDelimiter{\abs}{\lvert}{\rvert}
\providecommand{\norm}{}\let\norm\relax\DeclarePairedDelimiter{\norm}{\lVert}{\rVert}
\usepackage[table]{xcolor} % [table] : fournit aussi \rowcolors (alternance de lignes)
\usepackage{pagecolor}
\usepackage{tikz}
\usetikzlibrary{arrows.meta,positioning,calc,shapes.geometric,shapes.misc,shapes.arrows,decorations.pathmorphing,decorations.pathreplacing,decorations.markings,patterns,shadows,mindmap,trees,fit,backgrounds,matrix,intersections,angles,quotes}
\usepackage{pgfplots}
\usepackage[version=4]{mhchem} % équations nucléaires \ce{^{235}_{92}U}
\usepackage[european,siunitx]{circuitikz} % schémas électriques
\pgfplotsset{compat=1.18}
\usepgfplotslibrary{fillbetween,groupplots} % aires entre courbes (calcul intégral)
\usepackage{enumitem}
\usepackage{fancyhdr}
% [most] charge toutes les bibliothèques tcolorbox (listings, minted, raster,
% documentation...) et gonfle inutilement la mémoire TeX sur un document long
% (~300 boîtes) : on ne charge que ce qui est réellement utilisé.
\usepackage[skins,breakable]{tcolorbox}
\usepackage{graphicx}
\usepackage{colortbl}
\usepackage{booktabs}
\usepackage{tabularx}
\usepackage{multirow}
\usepackage{diagbox} % case d'en-tête diagonale des tableaux à double entrée
% Colonnes extensibles centrée / gauche / droite (C, L, R), souvent utilisées par les modèles
\newcolumntype{C}{>{\centering\arraybackslash}X}
\newcolumntype{L}{>{\raggedright\arraybackslash}X}
\newcolumntype{R}{>{\raggedleft\arraybackslash}X}
\usepackage{array}
\usepackage{multicol}
\usepackage{siunitx}
% parse-numbers=false : accepte aussi \SI{2,5 \times 10^{-3}}{...}
\sisetup{locale=FR,output-decimal-marker={,},group-minimum-digits=4,parse-numbers=false}
\DeclareSIUnit\ohm{\ensuremath{\symup{\Omega}}} % U+2126 absent de la police
\DeclareSIUnit\division{div} % réglages d'oscilloscope (ms/div, V/div)
\DeclareSIUnit\tr{tr} % tours (vitesse de rotation, ex. \SI{1500}{\tr\per\minute})
\DeclareSIUnit\molar{M} % concentration molaire (ex. \SI{3}{\molar}, \SI{1}{\micro\molar})
\DeclareSIUnit\an{an} % année (le modèle confond parfois avec \year, un compteur TeX) — ex. \SI{5}{\kilo\meter\per\an}
\DeclareSIUnit\FCFA{FCFA} % franc CFA (ex. \SI{500}{\FCFA\per\kilo\gram})
\DeclareSIUnit\degreeBrix{\textdegree Brix} % teneur en sucre d'un sirop/jus (réfractométrie)
\DeclareSIUnit\month{mois} % le modèle utilise parfois \month, un compteur TeX, comme unité de durée
\DeclareSIUnit\microliter{\micro\liter} % le modèle écrit parfois \microliter en un seul token
\DeclareSIUnit\million{millions} % dénombrement (ex. \SI{15}{\million\per\mL} spermatozoïdes)
\usepackage{qrcode}
% \degree : commande fréquemment utilisée par les modèles pour un angle ou une
% température en dehors de \SI{}{} (non fournie par les paquets chargés).
\providecommand{\degree}{^{\circ}}
% \qed (fin de démonstration), utilisé par les modèles sans amsthm
\providecommand{\qed}{\hfill$\square$}
% \barre : double barre des valeurs interdites dans un tableau de variations (tkz-tab)
\providecommand{\barre}{\ensuremath{\|}}
% environnement proof (amsthm non chargé) utilisé par les modèles
\ifdefined\proof\else\newenvironment{proof}[1][Démonstration]{\par\noindent\textit{#1.}\ }{\qed\par}\fi
\usepackage{lastpage}
\usepackage{hyperref}
\hypersetup{hidelinks}

% ============================================================
% Palette « Pop » OnBuch+ — mêmes hex que la classe P (plus_ui.dart)
% ============================================================
\definecolor{poporange}{HTML}{F59321}
\definecolor{popdark}{HTML}{E07A0C}
\definecolor{popcream}{HTML}{FFF6E8}
\definecolor{popink}{HTML}{1C1714}
\definecolor{poppurple}{HTML}{7A5AE0}
\definecolor{popgreen}{HTML}{2E9E6A}
\definecolor{poppink}{HTML}{E0457B}
\definecolor{popblue}{HTML}{2F7DE1}
\definecolor{popgold}{HTML}{C98A12}
\definecolor{popmuted}{HTML}{8A7F76}
\definecolor{popline}{HTML}{EADFD2}
\definecolor{popgray}{HTML}{8A7F76}
\colorlet{popgrey}{popgray}
% Alias tolérants (le modèle nomme parfois les couleurs différemment)
\colorlet{popwhite}{white}
\colorlet{popblack}{popink}
\colorlet{popred}{poppink}
\colorlet{popyellow}{popgold}
% Teintes claires (fonds de tableaux/encadrés) — le modèle nomme parfois
% les couleurs avec un suffixe L (ex. popblueL) sans les avoir définies.
\colorlet{poporangeL}{poporange!20}
\colorlet{popdarkL}{popdark!20}
\colorlet{poppurpleL}{poppurple!20}
\colorlet{popgreenL}{popgreen!20}
\colorlet{poppinkL}{poppink!20}
\colorlet{popblueL}{popblue!20}
\colorlet{popgoldL}{popgold!20}
\colorlet{popredL}{popred!20}
\colorlet{popgrayL}{popgray!20}
% Teintes claires pour fonds de tableaux / zones
\colorlet{popblueL}{popblue!10!white}
\colorlet{poporangeL}{poporange!14!white}
\colorlet{popgreenL}{popgreen!12!white}
\colorlet{poppurpleL}{poppurple!12!white}
\colorlet{poppinkL}{poppink!12!white}
\colorlet{popgoldL}{popgold!14!white}

\pagecolor{popcream}

% ============================================================
% Typographie
% ============================================================
\defaultfontfeatures{Ligatures=TeX}
\setmainfont{PlusJakartaSans-Medium.ttf}[Path=../../../fonts/,BoldFont=PlusJakartaSans-Bold.ttf]
% Maths en sans-serif (Fira Math) avec chiffres et lettres droites de Plus
% Jakarta, pour que les formules se fondent dans le texte.
\usepackage[math-style=ISO,bold-style=ISO]{unicode-math}
\setmathfont{FiraMath-Regular.otf}[Path=../../../fonts/]
\setmathfont{PlusJakartaSans-Medium.ttf}[Path=../../../fonts/,range={up/num,up/{Latin,latin},\mathup}]
\usepackage{icomma} % virgule décimale sans espace en mode math
% Caractères mathématiques Unicode absents des polices de texte (Plus Jakarta,
% Archivo Black) : rendus par leur équivalent mathématique, en texte comme en
% formule — sinon ils s'affichent en carré vide (ex. « Arithmétique dans ℤ »).
\usepackage{newunicodechar}
\newunicodechar{ℝ}{\ensuremath{\mathbb{R}}}
\newunicodechar{ℤ}{\ensuremath{\mathbb{Z}}}
\newunicodechar{ℂ}{\ensuremath{\mathbb{C}}}
\newunicodechar{ℕ}{\ensuremath{\mathbb{N}}}
\newunicodechar{ℚ}{\ensuremath{\mathbb{Q}}}
\newunicodechar{𝔻}{\ensuremath{\mathbb{D}}}
\newunicodechar{α}{\ensuremath{\alpha}}
\newunicodechar{β}{\ensuremath{\beta}}
\newunicodechar{γ}{\ensuremath{\gamma}}
\newunicodechar{δ}{\ensuremath{\delta}}
\newunicodechar{ε}{\ensuremath{\varepsilon}}
\newunicodechar{θ}{\ensuremath{\theta}}
\newunicodechar{λ}{\ensuremath{\lambda}}
\newunicodechar{μ}{\ensuremath{\mu}}
\newunicodechar{σ}{\ensuremath{\sigma}}
\newunicodechar{τ}{\ensuremath{\tau}}
\newunicodechar{φ}{\ensuremath{\varphi}}
\newunicodechar{ψ}{\ensuremath{\psi}}
\newunicodechar{ω}{\ensuremath{\omega}}
\newunicodechar{Σ}{\ensuremath{\Sigma}}
% majuscules produites par \MakeUppercase dans les titres (α-aminés -> Α)
\newunicodechar{Α}{\ensuremath{\alpha}}
\newunicodechar{Β}{\ensuremath{\beta}}
\newunicodechar{Γ}{\ensuremath{\Gamma}}
\newunicodechar{Δ}{\ensuremath{\Delta}}
\newunicodechar{Ε}{\ensuremath{\varepsilon}}
\newunicodechar{Θ}{\ensuremath{\Theta}}
\newunicodechar{Λ}{\ensuremath{\Lambda}}
\newunicodechar{Μ}{\ensuremath{\mu}}
\newunicodechar{Τ}{\ensuremath{\tau}}
\newunicodechar{Φ}{\ensuremath{\Phi}}
\newunicodechar{Ψ}{\ensuremath{\Psi}}
\newunicodechar{Ω}{\ensuremath{\Omega}}
\newunicodechar{↦}{\ensuremath{\mapsto}}
\newunicodechar{∈}{\ensuremath{\in}}
\newunicodechar{∉}{\ensuremath{\notin}}
\newunicodechar{∧}{\ensuremath{\wedge}}
\newunicodechar{⇔}{\ensuremath{\Leftrightarrow}}
\newunicodechar{⟺}{\ensuremath{\Longleftrightarrow}}
\newunicodechar{⇒}{\ensuremath{\Rightarrow}}
\newunicodechar{⟹}{\ensuremath{\Longrightarrow}}
\newunicodechar{∘}{\ensuremath{\circ}}
\newunicodechar{∪}{\ensuremath{\cup}}
\newunicodechar{∩}{\ensuremath{\cap}}
\newunicodechar{≡}{\ensuremath{\equiv}}
\newunicodechar{⊂}{\ensuremath{\subset}}
\newunicodechar{⊥}{\ensuremath{\perp}}
\newunicodechar{∃}{\ensuremath{\exists}}
\newunicodechar{∀}{\ensuremath{\forall}}
\newunicodechar{∛}{\ensuremath{\sqrt[3]{\,}}}
\newunicodechar{∜}{\ensuremath{\sqrt[4]{\,}}}
\newunicodechar{⃗}{\ensuremath{{}^{\to}}}
\newunicodechar{ⁿ}{\ifmmode^{n}\else\textsuperscript{\ensuremath{n}}\fi}
\newunicodechar{ˣ}{\ifmmode^{x}\else\textsuperscript{\ensuremath{x}}\fi}
\newunicodechar{ᵇ}{\ifmmode^{b}\else\textsuperscript{\ensuremath{b}}\fi}
\newunicodechar{ʳ}{\ifmmode^{r}\else\textsuperscript{\ensuremath{r}}\fi}
\newunicodechar{ᵘ}{\ifmmode^{u}\else\textsuperscript{\ensuremath{u}}\fi}
\newunicodechar{ʸ}{\ifmmode^{y}\else\textsuperscript{\ensuremath{y}}\fi}
\newunicodechar{ᵃ}{\ifmmode^{a}\else\textsuperscript{\ensuremath{a}}\fi}
\newunicodechar{ᵉ}{\ifmmode^{e}\else\textsuperscript{\ensuremath{e}}\fi}
\newunicodechar{ᵏ}{\ifmmode^{k}\else\textsuperscript{\ensuremath{k}}\fi}
\newunicodechar{ᵖ}{\ifmmode^{p}\else\textsuperscript{\ensuremath{p}}\fi}
\newunicodechar{ᴺ}{\ifmmode^{N}\else\textsuperscript{\ensuremath{N}}\fi}
\newunicodechar{ᶜ}{\ifmmode^{c}\else\textsuperscript{\ensuremath{c}}\fi}
\newunicodechar{ʰ}{\ifmmode^{h}\else\textsuperscript{\ensuremath{h}}\fi}
\newunicodechar{ᵠ}{\ifmmode^{q}\else\textsuperscript{\ensuremath{q}}\fi}
\newunicodechar{ₙ}{\ifmmode_{n}\else\textsubscript{\ensuremath{n}}\fi}
\newunicodechar{ₐ}{\ifmmode_{a}\else\textsubscript{\ensuremath{a}}\fi}
\newunicodechar{ᵢ}{\ifmmode_{i}\else\textsubscript{\ensuremath{i}}\fi}
\newunicodechar{ᵣ}{\ifmmode_{r}\else\textsubscript{\ensuremath{r}}\fi}
\newunicodechar{ₖ}{\ifmmode_{k}\else\textsubscript{\ensuremath{k}}\fi}
\newunicodechar{ᵦ}{\ifmmode_{\beta}\else\textsubscript{\ensuremath{\beta}}\fi}
\newunicodechar{ₓ}{\ifmmode_{x}\else\textsubscript{\ensuremath{x}}\fi}
\newfontfamily\popdisplay{ArchivoBlack-Regular.ttf}[Path=../../../fonts/]
\newfontfamily\poplogo{SpaceGrotesk-Bold.ttf}[Path=../../../fonts/]
\newfontfamily\popsemi{PlusJakartaSans-SemiBold.ttf}[Path=../../../fonts/]
\newfontfamily\popxbold{PlusJakartaSans-ExtraBold.ttf}[Path=../../../fonts/]
\newfontfamily\popxboldls{PlusJakartaSans-ExtraBold.ttf}[Path=../../../fonts/,LetterSpace=3.0]

\setlist[enumerate]{leftmargin=1.7em,itemsep=5pt,topsep=4pt}
\setlist[itemize]{leftmargin=1.7em,itemsep=4pt,topsep=4pt,label={\color{poporange}\textbullet}}
\setlength{\parindent}{0pt}
\setlength{\parskip}{5pt}
\renewcommand{\arraystretch}{1.25}

% pgfplots : style Pop par défaut
\pgfplotsset{
  every axis/.append style={axis line style={popink,line width=1pt},tick style={popink},
    grid style={popline,line width=0.5pt},label style={font=\small},tick label style={font=\footnotesize}},
  popcurve/.style={poporange,line width=1.8pt,no markers,smooth},
}
% Même style disponible dans un \draw TikZ simple (hors pgfplots)
\tikzset{popcurve/.style={poporange,line width=1.8pt}}
\tikzset{popgrid/.style={popline,very thin}}
\colorlet{popcurve}{poporange} % le nom du style est parfois utilisé comme couleur

% ============================================================
% Pill / squiggle
% ============================================================
\newcommand{\poppill}[3]{%
  \tikz[baseline=-0.55ex]\node[draw=popink,line width=1.2pt,rounded corners=2.6pt,%
    fill=#2,inner xsep=6.5pt,inner ysep=3pt]{%
    \popxboldls\fontsize{7}{7}\selectfont\color{#3}\MakeUppercase{#1}};%
}
\newcommand{\popsquiggle}[1]{%
  \tikz[baseline=-0.4ex]{
    \draw[#1,line width=2.6pt,line cap=round]
      plot [smooth] coordinates {(0,0.09) (0.34,0.01) (0.68,0.21) (1.02,0.01) (1.36,0.10)};
  }%
}
% Mot-clé surligné (marqueur orange sous le texte)
\newcommand{\cle}[1]{{\popxbold\color{popdark}#1}}

% ============================================================
% En-tête / pied
% ============================================================
\pagestyle{fancy}
\fancyhf{}
\renewcommand{\headrulewidth}{0pt}
\renewcommand{\footrulewidth}{0pt}
\lhead{\raisebox{-0.15ex}{\poplogo\fontsize{13}{13}\selectfont\color{popink}OnBuch\color{poporange}+}}
\rhead{\poppill{\DOCMATIERE\ \textbullet\ \DOCNIVEAU\ \textbullet\ Leçon \DOCLECON}{white}{popink}}
\lfoot{\popxbold\fontsize{7.6}{7.6}\selectfont\color{popmuted}\MakeUppercase{Cours signé OnBuch+}\ \textbullet\ {\popsemi\DOCTITRE}}
\rfoot{\tikz[baseline=-0.6ex]\node[rounded corners=8.5pt,fill=popink,minimum height=17pt,minimum width=17pt,inner xsep=5pt,inner sep=0]{\popxbold\fontsize{8}{8}\selectfont\color{popcream}\thepage\,/\,\pageref*{LastPage}};}

% ============================================================
% Titres
% ============================================================
\newcommand{\chaptitre}[2]{%
  \begin{tcolorbox}[enhanced,colback=poporange,colframe=popink,boxrule=2.6pt,arc=11pt,
      drop shadow=popink,left=15pt,right=15pt,top=12pt,bottom=11pt,before skip=3pt,after skip=18pt]
    \poppill{#2}{popink}{popcream}\par\vspace{9pt}
    {\raggedright\hyphenpenalty=10000\exhyphenpenalty=10000\popdisplay\fontsize{22}{24}\selectfont\color{popcream}\MakeUppercase{#1}\par}
  \end{tcolorbox}
}
% Section numérotée : pastille orange avec le numéro + titre Archivo
\newcounter{popsec}
\newcommand{\coursec}[1]{%
  \stepcounter{popsec}\setcounter{popsub}{0}%
  \par\vspace{16pt}\noindent
  \begin{minipage}{\linewidth}
  \tikz[baseline=-0.7ex]\node[draw=popink,line width=1.6pt,fill=poporange,rounded corners=4pt,minimum size=22pt,inner sep=0]{\popdisplay\fontsize{12}{12}\selectfont\color{popcream}\thepopsec};%
  \hspace{8pt}{\popdisplay\fontsize{15}{17}\selectfont\color{popink}\MakeUppercase{#1}}\par
  \vspace{4pt}\noindent{\color{poporange}\rule{36pt}{3.6pt}}
  \end{minipage}\par\nopagebreak\vspace{8pt}
}
\newcounter{popsub}
\newcommand{\courssub}[1]{%
  \stepcounter{popsub}%
  \par\vspace{9pt}\noindent{\popxbold\fontsize{11.5}{13}\selectfont\color{popink}{\color{poporange}\thepopsec.\thepopsub}\hspace{6pt}#1}\par\nopagebreak\vspace{4pt}
}

% ============================================================
% Boîtes pédagogiques — même recette : fond blanc, bordure épaisse colorée,
% ombre dure encre, badge plein. Titre optionnel [#1].
% ============================================================
\tcbset{popbase/.style={breakable,enhanced,colback=white,boxrule=2.3pt,arc=9pt,
  shadow={3pt}{-3pt}{0pt}{popink},left=14pt,right=14pt,top=11pt,bottom=11pt,before skip=11pt,after skip=15pt,
  fontupper=\popsemi\fontsize{10.6}{14.5}\selectfont\color{popink},
  fontlower=\popsemi\fontsize{10.6}{14.5}\selectfont\color{popink}}}
% Coche dessinée (le glyphe ✓ n'existe pas dans les polices du document)
\newcommand{\popcheck}[1]{\tikz[baseline=-0.5ex]\draw[#1,line width=1.6pt,line cap=round,line join=round](-0.13,0.0)--(-0.04,-0.09)--(0.13,0.1);}
\providecommand{\coche}{\popcheck{popgreen}}
\providecommand{\resultat}[1]{\par\smallskip\noindent\fcolorbox{popgreen}{popgreenL}{\parbox{\dimexpr\linewidth-2\fboxsep-2\fboxrule\relax}{\textbf{Résultat.} #1}}\par\smallskip}
\newcommand{\popboxhead}[3]{\poppill{#1}{#2}{white}\ifx\relax#3\relax\else\hspace{6pt}{\popxbold\fontsize{10.6}{12}\selectfont\color{popink}#3}\fi\par\vspace{7pt}}

\newtcolorbox{aretenir}[1][]{popbase,colframe=popgreen,before upper={\popboxhead{À retenir}{popgreen}{#1}}}
\newtcolorbox{exemplebox}[1][]{popbase,colframe=poporange,before upper={\popboxhead{Exemple}{poporange}{#1}}}
\newtcolorbox{definition}[1][]{popbase,colframe=popblue,before upper={\popboxhead{Définition}{popblue}{#1}}}
\newtcolorbox{propriete}[1][]{popbase,colframe=poppurple,before upper={\popboxhead{Propriété}{poppurple}{#1}}}
\newtcolorbox{methode}[1][]{popbase,colframe=popgold,before upper={\popboxhead{Méthode}{popgold}{#1}}}
\newtcolorbox{attention}[1][]{popbase,colframe=poppink,before upper={\popboxhead{Attention}{poppink}{#1}}}
\newtcolorbox{experience}[1][]{popbase,colframe=popblue,colback=popblueL,before upper={\popboxhead{Expérience}{popblue}{#1}}}
% « Le savais-tu ? » : carte violette pleine (contexte, culture, Cameroun)
\newtcolorbox{savaistu}[1][]{popbase,colback=poppurple,colframe=popink,
  fontupper=\popsemi\fontsize{10.4}{14.2}\selectfont\color{white},
  before upper={\poppill{Le savais-tu\,?}{white}{poppurple}\ifx\relax#1\relax\else\hspace{6pt}{\popxbold\color{white}#1}\fi\par\vspace{7pt}}}
% Exercice résolu : énoncé en haut, \tcblower puis la solution sur fond crème
\newcounter{popexo}\newcounter{popexr}
\newtcolorbox{exoresolu}[1][]{popbase,colframe=poporange,colbacklower=popcream,
  segmentation style={popink,line width=1.2pt,dashed},
  before upper={\stepcounter{popexr}\popboxhead{Exercice résolu \thepopexr}{poporange}{#1}},
  before lower={\poppill{Solution}{popink}{popcream}\par\vspace{6pt}}}
% Exercice (non résolu en place) — la correction est dans la partie Corrigés
\newtcolorbox{exercice}[1][]{popbase,colframe=popink,
  before upper={\stepcounter{popexo}\popboxhead{Exercice \thepopexo}{popink}{#1}}}
% Corrigé
\newtcolorbox{corrige}[1][]{popbase,colframe=popgreen,colback=popgreenL,
  before upper={\popboxhead{Corrigé}{popgreen}{#1}}}
% Objectifs / prérequis / bilan
\newtcolorbox{objectifs}{popbase,colframe=popink,colback=poporangeL,
  before upper={\popboxhead{Objectifs de la leçon}{popink}{}}}
\newtcolorbox{prerequis}{popbase,colframe=popink,colback=popblueL,
  before upper={\popboxhead{Prérequis}{popblue}{}}}
\newtcolorbox{bilan}{popbase,colframe=popink,colback=white,boxrule=2.8pt,
  before upper={\popboxhead{Fiche bilan}{poporange}{L'essentiel à connaître pour l'examen}}}
% Figure encadrée avec légende
\newtcolorbox{popfigure}[1][]{enhanced,breakable=false,colback=white,colframe=popline,boxrule=1.4pt,arc=8pt,
  left=10pt,right=10pt,top=10pt,bottom=8pt,before skip=10pt,after skip=12pt,halign=center,
  lower separated=false,halign lower=center,
  fontlower=\popsemi\fontsize{9}{11.5}\selectfont\color{popmuted},
  #1}
\newcommand{\legende}[1]{\par\vspace{6pt}{\popsemi\fontsize{9}{11.5}\selectfont\color{popmuted}#1}}

% ============================================================
% Signature OnBuch+
% ============================================================
\newcommand{\poplogoinline}[1]{{\poplogo\fontsize{#1}{#1}\selectfont\color{popink}OnBuch\color{poporange}+}}
\newcommand{\signatureonbuch}{%
  \begin{tcolorbox}[enhanced,colback=popink,colframe=popink,boxrule=2.2pt,arc=10pt,
      drop shadow=poporange,left=14pt,right=14pt,top=10pt,bottom=10pt,before skip=0pt,after skip=0pt]
    \tikz[baseline=-0.7ex]\node[circle,fill=popgreen,draw=popcream,line width=1.3pt,minimum size=22pt,inner sep=0]{\popcheck{white}};%
    \hspace{9pt}\begin{minipage}[c]{0.62\linewidth}
      {\popxbold\fontsize{12}{13}\selectfont\color{popcream}Signé\ \textbullet\ {\poplogo OnBuch\color{poporange}+}}\par\vspace{2pt}
      {\raggedright\popsemi\fontsize{8}{10.5}\selectfont\color{popline}\MakeUppercase{Équipe pédagogique OnBuch+}\\\MakeUppercase{Conforme au programme officiel MINESEC}\par}
    \end{minipage}\hfill
    {\poplogo\fontsize{22}{22}\selectfont\color{popcream}OnBuch\color{poporange}+}
  \end{tcolorbox}%
}

% ============================================================
% Page de garde
% #1 titre · #2 matière · #3 niveau · #4 module · #5 n° leçon · #6 durée ·
% #7 description / objectifs (texte ou itemize)
% ============================================================
\newcommand{\pagegarde}[7]{%
  \thispagestyle{empty}
  \newgeometry{a4paper,margin=1.7cm,top=1.6cm,bottom=1.4cm}%
  \begin{tikzpicture}[remember picture,overlay]
    \fill[poporange!18] ([xshift=-3.2cm,yshift=-3.6cm]current page.north east) circle (4.6cm);
    \fill[poppurple!14] ([xshift=2.4cm,yshift=7.6cm]current page.south west) circle (3.8cm);
  \end{tikzpicture}%
  {\poplogo\fontsize{30}{30}\selectfont\color{popink}OnBuch\color{poporange}+}\hfill
  \raisebox{0.6ex}{\poppill{Cours signé \textbullet\ Premium}{popink}{popcream}}\par
  \vspace{1pt}\popsquiggle{poppurple}\par
  \vspace{8pt}
  \begin{tcolorbox}[enhanced,colback=poporange,colframe=popink,boxrule=2.8pt,arc=13pt,
      drop shadow=popink,left=18pt,right=18pt,top=12pt,bottom=13pt,after skip=10pt]
    \poppill{#2 \textbullet\ #3}{popink}{popcream}\hspace{5pt}\poppill{#4}{white}{popink}\par\vspace{8pt}
    {\popdisplay\fontsize{14}{14}\selectfont\color{popink}LEÇON #5}\par\vspace{4pt}
    {\raggedright\hyphenpenalty=10000\exhyphenpenalty=10000\popdisplay\fontsize{25}{27}\selectfont\color{popcream}\MakeUppercase{#1}\par}
    \vspace{8pt}
    {\popxbold\fontsize{9.5}{9.5}\selectfont\color{popink}\MakeUppercase{Durée conseillée : #6}}
  \end{tcolorbox}
  \begin{tcolorbox}[enhanced,colback=white,colframe=popink,boxrule=2.2pt,arc=10pt,
      drop shadow=popink,left=16pt,right=16pt,top=10pt,bottom=10pt,after skip=8pt]
    \poppill{Dans cette leçon}{poppurple}{white}\par\vspace{5pt}
    \popsemi\fontsize{9.3}{12.6}\selectfont\color{popink}#7
  \end{tcolorbox}
  \noindent\begin{minipage}[t]{0.487\linewidth}
    \begin{tcolorbox}[enhanced,colback=popgreen,colframe=popink,boxrule=2.2pt,arc=10pt,
        drop shadow=popink,left=11pt,right=11pt,top=10pt,bottom=10pt,height=3.1cm,valign=center]
      \tcbox[colback=white,colframe=popink,boxrule=1.3pt,arc=5pt,boxsep=0pt,nobeforeafter,
        left=3pt,right=3pt,top=3pt,bottom=3pt,valign=center]{\qrcode[height=1.9cm]{https://play.google.com/store/apps/details?id=com.luvvix.onbuch&hl=fr}}%
      \hspace{8pt}\begin{minipage}[c]{0.5\linewidth}
        {\popdisplay\fontsize{11}{12.5}\selectfont\color{white}\MakeUppercase{Télécharge OnBuch}}\par\vspace{4pt}
        {\raggedright\popsemi\fontsize{8.4}{11}\selectfont\color{white}Gratuit \textbullet\ Google Play\\Tuteur IA, annales, concours, résultats\par}
      \end{minipage}
    \end{tcolorbox}
  \end{minipage}\hfill
  \begin{minipage}[t]{0.487\linewidth}
    \begin{tcolorbox}[enhanced,colback=poppurple,colframe=popink,boxrule=2.2pt,arc=10pt,
        drop shadow=popink,left=11pt,right=11pt,top=10pt,bottom=10pt,height=3.1cm,valign=center]
      \tikz[baseline=-0.7ex]\node[circle,fill=white,draw=popink,line width=1.3pt,minimum size=1.6cm,inner sep=0]{\popdisplay\fontsize{24}{24}\selectfont\color{poppurple}+};%
      \hspace{8pt}\begin{minipage}[c]{0.6\linewidth}
        {\popdisplay\fontsize{11}{12.5}\selectfont\color{white}\MakeUppercase{Passe à OnBuch+}}\par\vspace{4pt}
        {\raggedright\popsemi\fontsize{8.4}{11}\selectfont\color{white}Tous les cours signés, Tuteur IA illimité, fiches PDF — onglet OnBuch+ de l'app\par}
      \end{minipage}
    \end{tcolorbox}
  \end{minipage}
  \vspace{8pt}
  \popcontenu
  \vfill
  \signatureonbuch
  \restoregeometry
  \newpage
}

% Rangée « Ce cours contient » (page de garde)
\newcommand{\poptile}[3]{% #1 couleur #2 gros texte #3 légende
  \begin{tcolorbox}[enhanced,colback=#1,colframe=popink,boxrule=2pt,arc=9pt,shadow={2.5pt}{-2.5pt}{0pt}{popink},
      left=6pt,right=6pt,top=9pt,bottom=9pt,halign=center,valign=center,height=1.75cm,width=\linewidth,nobeforeafter]
    {\popdisplay\fontsize{12.5}{13}\selectfont\color{white}\MakeUppercase{#2}}\par\vspace{3pt}
    {\centering\popsemi\fontsize{7.8}{9.5}\selectfont\color{white}#3\par}
  \end{tcolorbox}}
\newcommand{\popcontenu}{%
  \par\noindent{\popxboldls\fontsize{7.5}{7.5}\selectfont\color{popmuted}CE COURS SIGNÉ CONTIENT}\par\vspace{7pt}
  \noindent\begin{minipage}[t]{0.235\linewidth}\poptile{popblue}{Cours}{complet, illustré, pas à pas}\end{minipage}\hfill
  \begin{minipage}[t]{0.235\linewidth}\poptile{poporange}{Méthodes}{et exercices résolus}\end{minipage}\hfill
  \begin{minipage}[t]{0.235\linewidth}\poptile{poppink}{Exercices}{type examen + corrigés}\end{minipage}\hfill
  \begin{minipage}[t]{0.235\linewidth}\poptile{popgreen}{Fiche bilan}{l'essentiel à retenir}\end{minipage}\par}

% Sommaire
\newtcolorbox{sommaire}{enhanced,colback=white,colframe=popink,boxrule=2.2pt,arc=10pt,
  drop shadow=popink,left=18pt,right=18pt,top=13pt,bottom=13pt,before skip=6pt,after skip=20pt,
  before upper={\poppill{Sommaire}{poppurple}{white}\par\vspace{9pt}\popsemi\fontsize{10.6}{14.5}\selectfont\color{popink}}}

% Signature de fin de cours — poussée tout en bas de la dernière page
\newcommand{\findecours}{%
  \par\vfill\nopagebreak
  \begin{center}\popsquiggle{poporange}\end{center}\vspace{4pt}
  \signatureonbuch
}
\AtBeginDocument{\def\llbracket{\mathopen{[\mkern-3.5mu[}}\def\rrbracket{\mathclose{]\mkern-3.5mu]}}} % intervalles d'entiers (glyphe ⟦ absent de la police)
% Glyphes absents de Fira Math : redéfinis à partir de glyphes disponibles
\AtBeginDocument{%
  \def\setminus{\mathbin{\backslash}}\let\smallsetminus\setminus
  \def\vdots{\mathord{\vbox{\baselineskip3.2pt\lineskiplimit0pt\kern4pt\hbox{.}\hbox{.}\hbox{.}}}}%
  \def\ddots{\mathinner{\mkern1mu\raise6.5pt\hbox{.}\mkern2mu\raise3.6pt\hbox{.}\mkern2mu\raise0.7pt\hbox{.}\mkern1mu}}%
  \def\triangle{\mathord{\scalebox{0.9}{$\symup{\Delta}$}}}%
  \def\nabla{\mathord{\raisebox{\depth}{\scalebox{1}[-1]{$\symup{\Delta}$}}}}%
  \def\checkmark{\popcheck{popdark}}%
  \def\star{\mathbin{\ast}}\def\bigstar{\mathord{\ast}}%
  \def\diamond{\mathbin{\scalebox{0.6}{$\Diamond$}}}%
  \def\bigcup{\mathop{\mathchoice{\raisebox{-0.3em}{\scalebox{1.7}{$\cup$}}}{\scalebox{1.25}{$\cup$}}{\cup}{\cup}}\displaylimits}%
  \def\bigcap{\mathop{\mathchoice{\raisebox{-0.3em}{\scalebox{1.7}{$\cap$}}}{\scalebox{1.25}{$\cap$}}{\cap}{\cap}}\displaylimits}%
  \def\lesseqqgtr{\mathrel{\lessgtr}}\def\bigoplus{\mathop{\oplus}\displaylimits}%
}
\providecommand{\ssi}{si et seulement si} % abréviation fréquente
\AtBeginDocument{\providecommand{\wideparen}{\overparen}} % arc de cercle

%% FIGFIX-BEGIN v5 — correctifs globaux des figures (docs/onbuch-generation/tools/figfix)
\usepackage{adjustbox}\usepackage{etoolbox}\tcbuselibrary{raster}
% toute figure TikZ/pgfplots plus large que la ligne est réduite à la largeur disponible
\BeforeBeginEnvironment{tikzpicture}{\begin{adjustbox}{max width=\linewidth}}
\AfterEndEnvironment{tikzpicture}{\end{adjustbox}}
% légendes pgfplots lisibles par-dessus les courbes
\pgfplotsset{every axis/.append style={legend style={fill=white,fill opacity=0.92,text opacity=1,draw=black!15,inner sep=3pt}}}
% étiquettes d'arêtes (syntaxe edge["..."]) avec fond blanc
\tikzset{every edge quotes/.append style={fill=white,inner sep=1.2pt}}
%% FIGFIX-END
\begin{document}

\pagegarde{Media and Communication}{Anglais}{1re Tech}{Anglais – classe de Première technique}{N05}{13 séances (fiche de progression harmonisée)}{This lesson equips Première technique students with the language skills needed to exchange information about verbal and non-verbal communication through TV, radio and print media, to use new Information and Communication Technologies, and to describe a workshop visit. It combines grammar (active/passive voices, used for/in + V-ing, preferences, WH-questions, punctuation), vocabulary, listening, reading and writing tasks anchored in the Cameroonian context.
\vspace{4pt}
\begin{multicols}{2}\small
\begin{itemize}
\item À la fin de cette leçon, je sais échanger des informations sur l'exploitation des informations verbales et non verbales à travers la TV, la radio et la presse écrite.
\item À la fin de cette leçon, je sais utiliser correctement les voix active et passive pour décrire des procédés techniques.
\item À la fin de cette leçon, je sais employer 'used for/in + verb-ing' et 'used with + noun' pour expliquer la fonction des TIC (www, MP3, USB flash drive, WhatsApp, Viber).
\item À la fin de cette leçon, je sais comprendre des textes oraux portant sur les nouvelles technologies de l'information et de la communication.
\item À la fin de cette leçon, je sais lire et exploiter un texte relatant la visite d'un atelier technique.
\item À la fin de cette leçon, je sais exprimer mes préférences (prefer, would rather, like... better) dans des situations de communication.
\end{itemize}\end{multicols}}

\begin{sommaire}
\begin{multicols}{2}
\begin{enumerate}
\item Speaking: Exchanging Information through the Media
\item Grammar Focus 1: Active and Passive Voices
\item Listening \& Vocabulary: New Information and Communication Technologies
\item Reading: Visiting a Workshop
\item Grammar Focus 2: Punctuation, WH- and How Questions
\item Writing \& Integration: Reporting a Workshop Visit
\item Activité d'intégration
\item Méthodes et exercices résolus
\item Exercices
\item Corrigés des exercices
\item Fiche bilan
\end{enumerate}
\end{multicols}
\end{sommaire}

\begin{objectifs}
\begin{itemize}
\item À la fin de cette leçon, je sais échanger des informations sur l'exploitation des informations verbales et non verbales à travers la TV, la radio et la presse écrite.
\item À la fin de cette leçon, je sais utiliser correctement les voix active et passive pour décrire des procédés techniques.
\item À la fin de cette leçon, je sais employer 'used for/in + verb-ing' et 'used with + noun' pour expliquer la fonction des TIC (www, MP3, USB flash drive, WhatsApp, Viber).
\item À la fin de cette leçon, je sais comprendre des textes oraux portant sur les nouvelles technologies de l'information et de la communication.
\item À la fin de cette leçon, je sais lire et exploiter un texte relatant la visite d'un atelier technique.
\item À la fin de cette leçon, je sais exprimer mes préférences (prefer, would rather, like... better) dans des situations de communication.
\item À la fin de cette leçon, je sais ponctuer correctement mes phrases et poser des questions avec WH- et How.
\item À la fin de cette leçon, je sais rédiger une composition sur la visite d'un atelier et la manipulation/maintenance de machines et gadgets.
\item À la fin de cette leçon, je sais appliquer ces notions à des situations concrètes de la vie courante camerounaise et justifier mes réponses.
\end{itemize}
\end{objectifs}

\begin{prerequis}
\begin{itemize}
\item Les temps de base de l'anglais (present simple, past simple, present continuous) vus en classe de Seconde.
\item Le vocabulaire élémentaire des médias (television, radio, newspaper) acquis au collège.
\item La formation des questions simples avec les auxiliaires do/does/did.
\item Les structures de base de la phrase anglaise (sujet – verbe – complément).
\item L'usage courant d'un téléphone mobile et d'applications de messagerie dans la vie quotidienne.
\end{itemize}
\end{prerequis}

\newpage

\coursec{Speaking: Exchanging Information through the Media}

\textbf{Situation d'accroche.} Every morning in Yaoundé, millions of Cameroonians switch on \cle{CRTV} or \cle{Équinoxe TV}, listen to the radio in taxis, read \emph{Cameroon Tribune}, or chat on WhatsApp and Facebook before breakfast. Yet how many of them can tell the difference between reliable information and a rumour, or explain in English how a phone-repair workshop works, like the ones in the Melen neighbourhood? How are the media and the new technologies changing the way we communicate, learn and work? This lesson gives you the language tools to \cle{exchange information}, express your \cle{preferences} and describe technical equipment in English.

\textbf{Problématique.} How can we use English to exchange verbal and non-verbal information through the TV, the radio and the print media?

\courssub{What are the media?}

Imagine you want to know the result of a football match between Canon of Yaoundé and Coton Sport of Garoua. You can watch the match report on television, listen to the commentary on the radio, or read the article in a newspaper the next morning. In each case, the message travels from one sender to millions of receivers at the same time. These channels of mass communication are called the \cle{media}.

\begin{definition}[Media]
The \cle{media} (singular: \emph{medium}) are the channels of mass communication that carry information to a large number of people at the same time. The three main traditional media are:
\begin{itemize}
\item \textbf{Television (TV)}: it combines moving images, sounds and spoken words.
\item \textbf{Radio}: it uses only sounds and spoken words.
\item \textbf{Print media}: newspapers and magazines, which use written words and photographs.
\end{itemize}
\end{definition}

In Cameroon, you already know many media houses. For television, there are \textbf{CRTV} (the public broadcaster), \textbf{Canal 2 International}, \textbf{STV} and \textbf{Équinoxe TV}. For radio, there are \textbf{CRTV Radio}, \textbf{Radio Balafon} and \textbf{Sweet FM}. For the print media, the most famous newspapers are \textbf{Cameroon Tribune} (bilingual), \textbf{Le Messager} and \textbf{Mutations}.

\begin{popfigure}
\begin{tikzpicture}[font=\small]
% TV column
\fill[popblueL] (0,0) rectangle (3.6,4.6);
\draw[popblue, thick] (0,0) rectangle (3.6,4.6);
\node[font=\bfseries, text=popblue] at (1.8,4.2) {TELEVISION};
\draw[popdark, thick, rounded corners] (0.7,2.7) rectangle (2.9,3.7);
\draw[popdark] (1.4,2.7) -- (1.2,2.4); \draw[popdark] (2.2,2.7) -- (2.4,2.4);
\node[text=popdark] at (1.8,3.2) {\emph{images + sound}};
\node[align=center] at (1.8,1.6) {CRTV\\ Canal 2 International\\ STV\\ Équinoxe TV};
% Radio column
\fill[poporangeL] (4.1,0) rectangle (7.7,4.6);
\draw[poporange, thick] (4.1,0) rectangle (7.7,4.6);
\node[font=\bfseries, text=poporange] at (5.9,4.2) {RADIO};
\draw[popdark, thick, rounded corners] (4.8,2.7) rectangle (7.0,3.7);
\draw[popdark] (6.6,3.7) -- (7.3,4.3);
\draw[popdark] (5.1,3.2) circle (0.18);
\node[text=popdark] at (6.0,3.2) {\emph{sound only}};
\node[align=center] at (5.9,1.6) {CRTV Radio\\ Radio Balafon\\ Sweet FM};
% Print column
\fill[popgreenL] (8.2,0) rectangle (11.8,4.6);
\draw[popgreen, thick] (8.2,0) rectangle (11.8,4.6);
\node[font=\bfseries, text=popgreen!60!black] at (10.0,4.2) {PRINT MEDIA};
\draw[popdark, thick] (8.9,2.7) rectangle (11.1,3.7);
\draw[popdark] (9.1,3.45) -- (10.9,3.45);
\draw[popdark] (9.1,3.2) -- (10.5,3.2);
\draw[popdark] (9.1,2.95) -- (10.7,2.95);
\node[align=center] at (10.0,1.6) {Cameroon Tribune\\ Le Messager\\ Mutations};
\end{tikzpicture}
\legende{The three main types of media in Cameroon, with examples of media houses. Television combines images and sound; radio uses sound only; print media use written words and photographs.}
\end{popfigure}

\begin{savaistu}[Did you know?]
Since the liberalisation of the audiovisual sector in Cameroon in the year 2000, more than \textbf{100 private radio stations} have been created across the country. Community radios now broadcast in English, French and many national languages, bringing information to villages that have no television signal.
\end{savaistu}

\courssub{Verbal and Non-verbal Information}

When you watch the news on CRTV, you receive two kinds of messages at the same time. The \cle{news anchor} (le présentateur) speaks words: this is \cle{verbal information}. But behind him, you can see images of a flooded road in Douala, the colours of a political party, or the logo of a ministry: this is \cle{non-verbal information}.

\begin{definition}[Verbal and non-verbal information]
\begin{itemize}
\item \cle{Verbal information} is a message conveyed by words, either \textbf{spoken} (a radio interview, a TV report) or \textbf{written} (a newspaper article, a headline).
\item \cle{Non-verbal information} is a message conveyed \textbf{without words}: pictures, photographs, gestures, facial expressions, sounds, music, colours, symbols and logos.
\end{itemize}
\end{definition}

Let us observe a concrete case. During a TV report about a school fair in Bafoussam, the journalist says: ``More than five thousand students attended the fair.'' That is verbal information. At the same time, the camera shows crowded stands, smiling students and a banner with the words ``Education for All''. The images confirm and complete the words: that is non-verbal information. A good communicator always pays attention to \textbf{both}, because sometimes the picture tells more than the words.

\begin{center}
\begin{tabularx}{\linewidth}{|l|X|X|}
\hline
\rowcolor{popblueL} \textbf{Criterion} & \textbf{Verbal information} & \textbf{Non-verbal information} \\
\hline
Channel & Words (spoken or written) & Images, gestures, sounds, colours, logos \\
\hline
TV example & The news anchor's words & Background images behind the anchor \\
\hline
Radio example & The journalist's report & Music, tone of voice, sound effects \\
\hline
Newspaper example & The article and the headline & The photograph and the layout \\
\hline
Advantage & Precise and detailed & Fast, emotional, understood by all \\
\hline
Risk & Can be long or boring & Can be ambiguous or manipulative \\
\hline
\end{tabularx}
\end{center}

\begin{attention}[A very frequent mistake]
In English, the word \cle{information} is \textbf{uncountable}. You can never say ``informations'' or ``an information''. Say instead: \emph{a piece of information}, \emph{some information}, \emph{an item of news}. Correct: ``The radio gave us \textbf{some useful information}.'' Wrong: ``The radio gave us \textbf{informations}.'' The same rule applies to \emph{news} (always singular: ``The news \textbf{is} good'') and \emph{equipment} (never ``equipments'').
\end{attention}

\courssub{Useful Expressions for Exchanging Information}

To exchange information from the media, English speakers use reporting expressions. Learn them by heart; they will be very useful at the Probatoire and the Baccalauréat technique.

\begin{aretenir}[Reporting expressions]
\begin{itemize}
\item \textbf{According to the news...} the match was postponed. (\emph{selon les informations...})
\item \textbf{I heard on the radio that...} the price of fuel will change.
\item \textbf{I read in Cameroon Tribune that...} a new technical college will open.
\item \textbf{I saw on TV that...} the Indomitable Lions won their match.
\item \textbf{The picture shows that...} the bridge is damaged.
\item \textbf{He said (that)...} the workshop was well equipped. (reported speech)
\end{itemize}
\end{aretenir}

\begin{methode}[How to conduct an information exchange]
Follow these four steps, for example when you interview a witness of a local event:
\begin{enumerate}
\item \textbf{Greet} the person politely: \emph{Good morning, thank you for accepting this interview.}
\item \textbf{Ask a WH-question}: \emph{What did you see? Where did it happen? When did the match start?}
\item \textbf{Listen} carefully and note the key words (verbal and non-verbal details).
\item \textbf{Report the answer} with reported speech: \emph{He said that the match started at 4 p.m.} / \emph{She told me that the fair was very crowded.}
\end{enumerate}
\end{methode}

\begin{exemplebox}[A short dialogue at a school fair]
\textbf{Journalist (A):} Good morning! I am reporting for the school magazine. What is happening here today?\\
\textbf{Witness (B):} Good morning. It is our annual science and technology fair.\\
\textbf{A:} According to the programme, how many schools are taking part?\\
\textbf{B:} I heard on the radio that twelve schools registered, and the picture on the poster shows that there are also workshops on phone repair.\\
\textbf{A:} Thank you. So, you said that twelve schools are taking part and that there are technical workshops.\\
\textbf{B:} Exactly!
\end{exemplebox}

\courssub{Practice: Role-play and Guided Debate}

\begin{exoresolu}[Role-play: the journalist and the witness]
\textbf{Task.} In pairs, one student is a journalist from Équinoxe TV, the other is a witness of a local event (a football match or a school fair). The journalist interviews the witness, then reports two pieces of information to the class.

\tcblower
\textbf{Model answer.}
\begin{itemize}
\item \textbf{Journalist:} Good afternoon. I am from Équinoxe TV. What did you see at the stadium yesterday?
\item \textbf{Witness:} I saw a very exciting match. Canon beat Coton Sport two goals to one.
\item \textbf{Journalist:} How did the supporters react?
\item \textbf{Witness:} They sang and waved green and yellow flags.
\item \textbf{Report to the class:} ``He said that Canon beat Coton Sport two goals to one, and he told me that the supporters waved green and yellow flags.''
\end{itemize}
\end{exoresolu}

\begin{exercice}[Guided debate: which medium is best for a technician?]
In groups of four, discuss this question: \emph{Which medium is the most useful for a Cameroonian technician: TV, radio or newspapers?} Give one advantage and one disadvantage of each medium. Use the expressions: \emph{In my opinion...}, \emph{I think that...}, \emph{According to me, the advantage of radio is that...}
\end{exercice}

\begin{corrige}[Guided debate]
Possible ideas (accept all reasonable answers in correct English):
\begin{itemize}
\item \textbf{TV:} advantage --- it shows images, so you can see how a machine works; disadvantage --- it is expensive and needs electricity.
\item \textbf{Radio:} advantage --- it is cheap and reaches villages; disadvantage --- there are no pictures, so technical demonstrations are difficult.
\item \textbf{Newspapers:} advantage --- you can keep the article and read it again; disadvantage --- they arrive late and many people in rural areas cannot buy them daily.
\end{itemize}
\end{corrige}

\begin{exercice}[Verbal or non-verbal?]
Classify the following items: (1) the headline of \emph{Mutations}; (2) the logo of CRTV; (3) the journalist's commentary on Radio Balafon; (4) the photograph of a car accident in \emph{Le Messager}; (5) the music at the beginning of the news.
\end{exercice}

\begin{corrige}[Verbal or non-verbal?]
(1) Verbal (written words). (2) Non-verbal (symbol). (3) Verbal (spoken words). (4) Non-verbal (photograph). (5) Non-verbal (sound).
\end{corrige}

\begin{exercice}[Report the information]
Rewrite each sentence with the reporting expression in brackets. (1) ``The bridge in Wouri is damaged.'' (\emph{According to the news...}) (2) ``A new technical college will open in Bafoussam.'' (\emph{I read in Cameroon Tribune that...}) (3) ``The workshop is well equipped,'' said the visitor. (\emph{He said that...})
\end{exercice}

\begin{corrige}[Report the information]
(1) According to the news, the bridge in Wouri is damaged. (2) I read in Cameroon Tribune that a new technical college will open in Bafoussam. (3) He said that the workshop was well equipped. (Note the tense change in reported speech: \emph{is} becomes \emph{was}.)
\end{corrige}

\begin{aretenir}[Key points to remember]
\begin{itemize}
\item The \cle{media} are channels of mass communication: TV, radio and print media.
\item \cle{Verbal information} uses words; \cle{non-verbal information} uses images, gestures, sounds, colours and logos.
\item To report information, use: \emph{According to the news...}, \emph{I heard on the radio that...}, \emph{He said that...}.
\item \emph{Information} is uncountable: never ``informations''.
\end{itemize}
\end{aretenir}

\coursec{Grammar Focus 1: Active and Passive Voices}

In the previous section, you learned to exchange information through the media: you spoke about TV programmes, radio shows and newspapers. You used sentences like \emph{``The journalist writes the article''} and \emph{``CRTV broadcasts the news''}. In all these sentences, the person who \emph{does} the action is the subject. But in the world of media and technology, we often care more about the \emph{result} of the action than about the person who does it. When you read \emph{``The news is broadcast at 8 p.m.''}, the important thing is the news, not the person who reads it. This new way of building a sentence is called the \cle{passive voice}, and it is the grammar focus of this section.

\courssub{The Active Voice: A Quick Reminder}

\begin{definition}[Active voice]
In the \cle{active voice}, the \textbf{subject performs the action}. The basic structure is:
\begin{center}
\textbf{Subject + Verb + Object}
\end{center}
The subject is the \emph{agent} (le \emph{faiseur} de l'action), the verb shows the action, and the object receives the action.
\end{definition}

\begin{exemplebox}[Active sentences in the media world]
\begin{itemize}
\item \textbf{The journalist} \emph{writes} \textbf{the article}. (Le journaliste écrit l'article.)
\item \textbf{CRTV} \emph{broadcasts} \textbf{the news} at 8 p.m.
\item \textbf{The technician} \emph{repairs} \textbf{the machines} every week.
\item \textbf{The editor} \emph{has published} \textbf{the magazine}.
\item \textbf{The camera operator} \emph{films} \textbf{the report}. (L'opérateur filme le reportage.)
\end{itemize}
In each sentence, the subject (in bold at the beginning) \emph{does} the action.
\end{exemplebox}

Now, imagine a different situation. You arrive at the workshop and you see that the machines are working again. You do not know who repaired them, and honestly, it does not matter. What matters is the machines. In English, you cannot say \emph{``Someone repaired the machines''} every time; instead, you put the machines in front: \emph{``The machines are repaired.''} Welcome to the passive voice.

\courssub{The Passive Voice: Definition and Formation}

\begin{definition}[Passive voice]
In the \cle{passive voice}, the \textbf{subject receives the action instead of performing it}. The object of the active sentence becomes the new subject. The structure is:
\begin{center}
\textbf{Object (new subject) + be (at the tense of the active verb) + past participle (+ by + agent)}
\end{center}
The agent (celui qui fait l'action) is introduced by \emph{by}, but it is often omitted.
\end{definition}

\begin{propriete}[Only transitive verbs can be passive]
Only \cle{transitive verbs} — verbs followed by a direct object (un COD) — can be used in the passive voice. If there is no object, there is nothing to put in front.
\begin{itemize}
\item \emph{The journalist writes \textbf{the article}.} $\rightarrow$ passive possible: \emph{The article is written.}
\item \emph{The presenter arrived late.} $\rightarrow$ passive \textbf{impossible}: \emph{arrive} has no object.
\item \emph{The signal travels by satellite.} $\rightarrow$ passive \textbf{impossible}: \emph{travel} is intransitive here.
\end{itemize}
\end{propriete}

\begin{popfigure}
\begin{center}
\begin{tikzpicture}[font=\small]
% Active row
\node[draw=popblue, fill=popblueL, rounded corners, minimum width=2.6cm, minimum height=0.9cm] (S) at (0,0) {\textbf{Subject}};
\node[draw=poporange, fill=poporangeL, rounded corners, minimum width=2.2cm, minimum height=0.9cm] (V) at (3.4,0) {\textbf{Verb}};
\node[draw=popgreen, fill=popgreenL, rounded corners, minimum width=2.6cm, minimum height=0.9cm] (O) at (6.4,0) {\textbf{Object}};
\node[above=0.05cm of S, popdark] {\scriptsize The journalist};
\node[below=0.05cm of V, popdark] {\scriptsize writes};
\node[below=0.05cm of O, popdark] {\scriptsize the article};
\node[popdark] at (3.2,1.05) {\textbf{ACTIVE}};
% Big transformation arrow
\draw[-{Stealth[length=4mm]}, line width=2.5pt, poppurple] (8.1,0) -- (9.6,0);
% Passive row
\node[draw=popgreen, fill=popgreenL, rounded corners, minimum width=2.6cm, minimum height=0.9cm] (O2) at (11.2,0) {\textbf{New Subject}};
\node[draw=poppink, fill=poppinkL, rounded corners, minimum width=1.6cm, minimum height=0.9cm] (BE) at (13.6,0) {\textbf{be}};
\node[draw=poporange, fill=poporangeL, rounded corners, minimum width=2.2cm, minimum height=0.9cm] (VPP) at (16.1,0) {\textbf{V--pp}};
\node[draw=popblue, fill=popblueL, rounded corners, minimum width=2.6cm, minimum height=0.9cm] (BY) at (19.2,0) {\textbf{(by + agent)}};
\node[above=0.05cm of O2, popdark] {\scriptsize The article};
\node[below=0.05cm of BE, popdark] {\scriptsize is};
\node[below=0.05cm of VPP, popdark] {\scriptsize written};
\node[below=0.05cm of BY, popdark] {\scriptsize by the journalist};
\node[popdark] at (15.2,1.05) {\textbf{PASSIVE}};
% Curved arrows showing movement
\draw[-{Stealth}, dashed, popgreen, thick] (O.north) .. controls (7.5,1.8) and (10,1.8) .. (O2.north);
\draw[-{Stealth}, dashed, popblue, thick] (S.south) .. controls (4,-1.6) and (16,-2.2) .. (BY.south);
\end{tikzpicture}
\end{center}
\legende{From active to passive: the \textbf{object} moves to the front and becomes the new subject; the verb becomes \emph{be + past participle}; the old subject moves to the end after \emph{by} (or disappears).}
\end{popfigure}

\begin{methode}[How to change an active sentence into a passive one]
Follow these four steps, in this order:
\begin{enumerate}
\item \textbf{Find the direct object} (le COD) of the active sentence. Ask: \emph{writes what? repairs what? films what?}
\item \textbf{Put this object at the beginning} of the new sentence: it becomes the subject.
\item \textbf{Conjugate the verb \emph{be}} at the same tense as the active verb, and \textbf{make it agree} with the new subject (singular or plural).
\item \textbf{Add the past participle} of the main verb, then \emph{by + agent} if the agent is important.
\end{enumerate}
Example: \emph{The technician repairs the machines.}
\begin{enumerate}
\item Object = \emph{the machines}.
\item \emph{The machines \ldots}
\item Active verb \emph{repairs} = present simple; new subject is plural $\rightarrow$ \emph{are}.
\item \emph{The machines \textbf{are repaired} (by the technician).}
\end{enumerate}
\end{methode}

\courssub{The Passive at the Main Tenses}

The golden rule: \textbf{only the verb \emph{be} changes}; the past participle never changes. The tense of the active verb is carried by \emph{be}. Study the table below with the verb \emph{to broadcast} (diffuser), a typical media verb.

\begin{center}
\begin{tabularx}{\linewidth}{|l|X|X|}
\hline
\rowcolor{popblueL} \textbf{Tense} & \textbf{Active} & \textbf{Passive} \\
\hline
Present simple & CRTV broadcasts the news. & The news \textbf{is broadcast} by CRTV. \\
\hline
Past simple & CRTV broadcast the match. & The match \textbf{was broadcast} by CRTV. \\
\hline
Present continuous & CRTV is broadcasting the debate. & The debate \textbf{is being broadcast}. \\
\hline
Past continuous & CRTV was broadcasting the show. & The show \textbf{was being broadcast}. \\
\hline
Future (will) & CRTV will broadcast the concert. & The concert \textbf{will be broadcast}. \\
\hline
Present perfect & CRTV has broadcast the report. & The report \textbf{has been broadcast}. \\
\hline
Past perfect & CRTV had broadcast the news. & The news \textbf{had been broadcast}. \\
\hline
Modals (can, must, should) & You \textbf{must repair} the generator. & The generator \textbf{must be repaired}. \\
\hline
\end{tabularx}
\end{center}

\begin{aretenir}[Key structures to memorize]
\begin{itemize}
\item Present simple: \textbf{am / is / are + V-pp}
\item Past simple: \textbf{was / were + V-pp}
\item Present continuous: \textbf{am / is / are + being + V-pp}
\item Past continuous: \textbf{was / were + being + V-pp}
\item Future: \textbf{will be + V-pp}
\item Present perfect: \textbf{have / has + been + V-pp}
\item Past perfect: \textbf{had + been + V-pp}
\item Modals: \textbf{modal + be + V-pp} (ex: \emph{must be done, can be sent, should be checked})
\end{itemize}
Note: \emph{broadcast} is an irregular verb: broadcast -- broadcast -- broadcast. Other useful media and workshop verbs: \emph{write -- wrote -- written}, \emph{send -- sent -- sent}, \emph{make -- made -- made}, \emph{repair -- repaired -- repaired}, \emph{transmit -- transmitted -- transmitted}, \emph{install -- installed -- installed}, \emph{record -- recorded -- recorded}, \emph{film -- filmed -- filmed}, \emph{edit -- edited -- edited}.
\end{aretenir}

\begin{attention}[Frequent error: agreement of \emph{be}]
The verb \emph{be} must agree with the \textbf{new subject} (the former object), not with the agent.
\begin{itemize}
\item Correct: \emph{The news \textbf{is} broadcast at 8 p.m.} (\emph{news} is singular in English.)
\item Wrong: \emph{The news \textbf{are} broadcast.}
\item Correct: \emph{The machines \textbf{are} repaired every week.}
\item Wrong: \emph{The machines \textbf{is} repaired.}
\end{itemize}
Another frequent error: forgetting \emph{be} completely. \emph{``The article written by the journalist''} is not a complete sentence — say \emph{``The article \textbf{is} written by the journalist.''}
\end{attention}

\courssub{When Do We Use the Passive Voice?}

\begin{propriete}[Three main uses of the passive]
We prefer the passive voice when the agent is:
\begin{enumerate}
\item \textbf{Unknown} (inconnu): \emph{My USB flash drive has been stolen!} (We do not know who stole it.)
\item \textbf{Obvious} (évident): \emph{The thief was arrested.} (Obviously by the police.)
\item \textbf{Unimportant} (sans importance): \emph{The machines are repaired every week.} (Who repairs them does not matter; the result matters.)
\end{enumerate}
We keep \emph{by + agent} only when the agent adds real information: \emph{The interview was conducted by our best journalist.}
\end{propriete}

\begin{savaistu}[The passive in technical and scientific English]
The passive voice is the favourite structure of technical documents, manuals and news reports, because it sounds objective and focuses on facts and processes. In a workshop manual you will read: \emph{``The signal is transmitted by satellite. The antenna is connected to the receiver. The oil is changed every 5\,000 km.''} In Cameroon, CRTV news bulletins, radio programmes on CRTV National Station, and newspapers like \emph{Cameroon Tribune} constantly use the passive: \emph{``A new transmitter was installed in Bamenda.''} Mastering the passive will help you read technical manuals in your professional life.
\end{savaistu}

\begin{exemplebox}[Describing a technical process with the passive (verbal \& non-verbal)]
How a radio programme reaches your phone or radio set:
\begin{enumerate}
\item The presenter's voice \textbf{is recorded} in the studio. (verbal)
\item The sound \textbf{is transformed} into an electrical signal. (non-verbal signal)
\item The signal \textbf{is transmitted} by satellite or by antenna.
\item The signal \textbf{is received} by your radio or smartphone.
\item The programme \textbf{is heard} by thousands of listeners.
\end{enumerate}
Notice: nobody mentions the technicians. In a process (un procédé), the steps matter more than the people.
\end{exemplebox}

\begin{exoresolu}[From active to passive: a media example]
\textbf{Statement.} Transform this sentence into the passive voice: \emph{``CRTV broadcasts 3 news bulletins daily.''}
\tcblower
\textbf{Detailed solution.}
\begin{enumerate}
\item \textbf{Find the object}: broadcasts \emph{what?} $\rightarrow$ \emph{3 news bulletins}.
\item \textbf{Put it first}: \emph{3 news bulletins \ldots}
\item \textbf{Conjugate \emph{be}}: the active verb \emph{broadcasts} is present simple; the new subject \emph{3 news bulletins} is plural $\rightarrow$ \emph{are}.
\item \textbf{Add the past participle}: broadcast $\rightarrow$ broadcast (irregular verb). The agent \emph{CRTV} is meaningful, so we keep it with \emph{by}.
\end{enumerate}
\textbf{Answer:} \emph{3 news bulletins \textbf{are broadcast} daily \textbf{by CRTV}.}
\end{exoresolu}

\begin{exoresolu}[From passive to active: a workshop example]
\textbf{Statement.} Put this sentence back into the active voice: \emph{``The broken antenna was repaired by the technician yesterday.''}
\tcblower
\textbf{Detailed solution.}
\begin{enumerate}
\item The agent after \emph{by} is \emph{the technician}: it becomes the subject of the active sentence.
\item The passive structure is \emph{was repaired} = past simple passive. So the active verb is at the past simple: \emph{repaired}.
\item The passive subject \emph{the broken antenna} becomes the object.
\end{enumerate}
\textbf{Answer:} \emph{The technician \textbf{repaired} the broken antenna yesterday.}
\end{exoresolu}

\courssub{Media Vocabulary: Key Terms (Verbal \& Non-Verbal)}

\begin{aretenir}[Essential vocabulary for this unit]
\begin{center}
\begin{tabularx}{\linewidth}{|l|X|l|}
\hline
\rowcolor{popblueL} \textbf{English Term} & \textbf{Definition / Context} & \textbf{French} \\
\hline
\cle{Broadcast} (v/n) & To send out programmes on TV/radio & Diffuser / Diffusion \\
\cle{Transmit} (v) & To send signals (radio, TV, data) & Transmettre \\
\cle{Signal} (n) & Electrical / electromagnetic wave carrying info & Signal \\
\cle{Transmitter} (n) & Equipment that sends out signals & Émetteur \\
\cle{Receiver} (n) & Device that picks up signals (radio, phone) & Récepteur \\
\cle{Antenna / Aerial} (n) & Metal structure for sending/receiving waves & Antenne \\
\cle{News bulletin} (n) & Short news programme & Journal parlé / Flash info \\
\cle{Report / Coverage} (n) & Detailed account of an event & Reportage / Couverture \\
\cle{Editor} (n) & Person responsible for content & Rédacteur en chef \\
\cle{Technician} (n) & Skilled worker maintaining equipment & Technicien \\
\cle{Generator} (n) & Machine producing electricity & Groupe électrogène \\
\cle{Flash drive / USB key} (n) & Portable storage device & Clé USB \\
\cle{Workshop} (n) & Place for repair / practical work & Atelier \\
\hline
\end{tabularx}
\end{center}
\end{aretenir}

\courssub{Practice: Transforming Sentences}

\begin{exercice}[Active $\rightarrow$ passive: media and workshop]
Put the following sentences into the passive voice. Keep the agent only if it is useful.
\begin{enumerate}
\item The journalist writes the article every morning.
\item The students recorded the interview yesterday.
\item The station will broadcast the football match on Saturday.
\item The engineer has installed the new transmitter.
\item The technicians are repairing the generator.
\item You must check the antenna connection.
\item The camera operator is filming the scene.
\end{enumerate}
\end{exercice}

\begin{corrige}[Exercise: Active $\rightarrow$ passive]
\begin{enumerate}
\item \emph{The article \textbf{is written} by the journalist every morning.} (present simple, singular $\rightarrow$ is)
\item \emph{The interview \textbf{was recorded} by the students yesterday.} (past simple $\rightarrow$ was)
\item \emph{The football match \textbf{will be broadcast} on Saturday.} (the agent is unimportant, we omit it)
\item \emph{The new transmitter \textbf{has been installed} by the engineer.} (present perfect $\rightarrow$ has been)
\item \emph{The generator \textbf{is being repaired} by the technicians.} (present continuous $\rightarrow$ is being)
\item \emph{The antenna connection \textbf{must be checked}.} (modal \emph{must} $\rightarrow$ must be + V-pp)
\item \emph{The scene \textbf{is being filmed}.} (present continuous, agent obvious/unimportant)
\end{enumerate}
\end{corrige}

\begin{exercice}[Passive $\rightarrow$ active]
Put the following sentences into the active voice.
\begin{enumerate}
\item The news is read by the presenter at 7 p.m.
\item The machines were cleaned by the apprentices.
\item A new radio programme has been created by the young journalists.
\item The signal can be received by any smartphone.
\item The oil should be changed every 5\,000 km.
\end{enumerate}
\end{exercice}

\begin{corrige}[Exercise: Passive $\rightarrow$ active]
\begin{enumerate}
\item \emph{The presenter \textbf{reads} the news at 7 p.m.}
\item \emph{The apprentices \textbf{cleaned} the machines.}
\item \emph{The young journalists \textbf{have created} a new radio programme.}
\item \emph{Any smartphone \textbf{can receive} the signal.}
\item \emph{You / The technician \textbf{should change} the oil every 5\,000 km.}
\end{enumerate}
\end{corrige}

\begin{aretenir}[What you must remember]
\begin{itemize}
\item Active: \textbf{Subject + Verb + Object} — the subject \emph{does} the action.
\item Passive: \textbf{Object + be (tense) + past participle (+ by + agent)} — the subject \emph{receives} the action.
\item Only transitive verbs can become passive.
\item Only \emph{be} changes with the tense; the past participle never changes.
\item Use the passive when the agent is unknown, obvious or unimportant — very common in news and technical descriptions: \emph{The signal is transmitted by satellite. The generator must be repaired.}
\end{itemize}
\end{aretenir}

You can now describe what the media and the machines \emph{do}, but also what \emph{is done} to them. In the next section, \emph{Listening \& Vocabulary: New Information and Communication Technologies}, you will hear texts about the www, MP3 players, USB flash drives, WhatsApp and Viber — and you will notice how often the passive voice appears when people describe these technologies.

\coursec{Listening \& Vocabulary: New Information and Communication Technologies}

In the previous section, you learned how to change a sentence from the \cle{active voice} to the \cle{passive voice}: \emph{``Millions of Cameroonians use WhatsApp''} becomes \emph{``WhatsApp is used by millions of Cameroonians.''} Notice something interesting in that passive sentence: it tells us about a \textbf{use}, a \textbf{function}. In this section, we go further. We will learn the vocabulary of the \cle{New Information and Communication Technologies} (ICTs), train our ears to understand spoken texts about them, and master two very useful grammar structures: \emph{used for + verb-ing} and \emph{used with + noun}.

\courssub{What Are ICTs?}

\begin{definition}[ICT — Information and Communication Technologies]
\cle{ICTs} are technologies used to \textbf{create}, \textbf{store}, \textbf{exchange} and \textbf{use} information in \textbf{digital form}. They include the Internet (www), computers, smartphones, and all the tools and applications we use to communicate: e-mail, WhatsApp, Viber, MP3 players, USB drives, etc.
\end{definition}

Think of your daily life in Yaoundé, Douala, Bamenda or Bafoussam. In the morning, you check your \cle{smartphone} for messages. At school, the computer teacher shows you a \cle{website}. In the evening, you \cle{download} music or \cle{chat} with friends. All these actions belong to the world of ICTs. They are not luxury objects anymore: they are tools for learning, working and communicating.

\begin{savaistu}[ICTs in Cameroon]
Cameroon has more than \textbf{10 million mobile Internet subscribers}, mainly through the operators \textbf{MTN} and \textbf{Orange}. Mobile money services (MTN Mobile Money, Orange Money) allow people to send and receive money without a bank account. For many young Cameroonians, the smartphone is the first — and sometimes the only — computer they use.
\end{savaistu}

\courssub{Vocabulary Building: The Language of ICTs}

Here is the essential vocabulary for this lesson. Read each word aloud, then learn it with its French meaning. Pay attention to the pronunciation hints in brackets.

\begin{center}
\begin{tabularx}{\linewidth}{|l|X|l|}
\hline
\rowcolor{popblueL} \textbf{English word} & \textbf{Meaning / example} & \textbf{French} \\
\hline
www (World Wide Web) & the system of linked pages on the Internet & la toile mondiale \\
\hline
website & a place on the Internet with pages of information & site web \\
\hline
MP3 player & a small device that plays music files & lecteur MP3 \\
\hline
USB flash drive / pen drive & a small device used to store and carry files & clé USB \\
\hline
smartphone & a mobile phone with Internet and applications & téléphone intelligent \\
\hline
e-mail & an electronic message sent through the Internet & courriel \\
\hline
WhatsApp & a free messaging application for text, voice and video & WhatsApp \\
\hline
Viber & an application for free calls and messages & Viber \\
\hline
download & to copy a file from the Internet to your device & télécharger (vers soi) \\
\hline
upload & to send a file from your device to the Internet & téléverser, envoyer \\
\hline
browse & to look at pages on the Internet & naviguer, surfer \\
\hline
chat & to exchange written messages in real time & clavarder, discuter \\
\hline
\end{tabularx}
\end{center}

\begin{attention}[Download or upload?]
Students often confuse these two words. Remember the direction of the arrow: \textbf{down}load = the file comes \textbf{down} to your phone; \textbf{up}load = the file goes \textbf{up} from your phone to the Internet. You \emph{download} a song; you \emph{upload} a photo to WhatsApp status.
\end{attention}

\begin{popfigure}
\begin{center}
\begin{tikzpicture}[scale=1.0]
\node[circle, draw=popdark, fill=popblue, text=white, font=\bfseries\large, minimum size=2cm] (ict) at (0,0) {ICTs};
\node[draw=popdark, fill=poporangeL, rounded corners, align=center, font=\small] (www) at (90:3.6) {\textbf{www / website}\\ browsing information};
\node[draw=popdark, fill=popgreenL, rounded corners, align=center, font=\small] (mp3) at (30:3.9) {\textbf{MP3 player}\\ listening to music};
\node[draw=popdark, fill=popgoldL, rounded corners, align=center, font=\small] (usb) at (-30:3.9) {\textbf{USB drive}\\ storing and carrying files};
\node[draw=popdark, fill=poppinkL, rounded corners, align=center, font=\small] (wa) at (-90:3.6) {\textbf{WhatsApp}\\ chatting and sending files};
\node[draw=popdark, fill=poppurpleL, rounded corners, align=center, font=\small] (vib) at (-150:3.9) {\textbf{Viber}\\ free calls and messages};
\node[draw=popdark, fill=popblueL, rounded corners, align=center, font=\small] (mail) at (150:3.9) {\textbf{e-mail}\\ formal written messages};
\draw[very thick, popline] (ict) -- (www);
\draw[very thick, popline] (ict) -- (mp3);
\draw[very thick, popline] (ict) -- (usb);
\draw[very thick, popline] (ict) -- (wa);
\draw[very thick, popline] (ict) -- (vib);
\draw[very thick, popline] (ict) -- (mail);
\end{tikzpicture}
\end{center}
\legende{The ICT galaxy: six common tools and their main functions. Each tool answers the question: \emph{What is it used for?}}
\end{popfigure}

\courssub{Pronunciation Corner: Key ICT Words}

\begin{center}
\begin{tabularx}{\linewidth}{|l|l|X|}
\hline
\rowcolor{popblueL} \textbf{Word} & \textbf{Phonetics (approx.)} & \textbf{Tip} \\
\hline
website & /ˈweb.saɪt/ & two syllables: \emph{web-site} \\
\hline
smartphone & /ˈsmɑːt.fəʊn/ & stress on \emph{smart} \\
\hline
download & /ˌdaʊnˈləʊd/ & stress on \emph{load} \\
\hline
upload & /ˌʌpˈləʊd/ & stress on \emph{load} \\
\hline
browse & /braʊz/ & rhymes with \emph{house} \\
\hline
chat & /tʃæt/ & like \emph{cat} with \emph{ch} \\
\hline
USB & /ˌjuː.esˈbiː/ & spell the letters: U-S-B \\
\hline
MP3 & /ˌem.piːˈθriː/ & say: M-P-three \\
\hline
\end{tabularx}
\end{center}

\courssub{Listening Skills: How to Understand a Spoken Text}

Listening is often the most difficult skill for students. The good news: you do not need to understand \emph{every} word. You need a \textbf{strategy}.

\begin{methode}[Three steps to succeed in a listening task]
\begin{enumerate}
\item \textbf{Before listening:} read the questions carefully. They tell you what to listen for. Anticipate: if the question is \emph{``How much does the bundle cost?''}, listen for numbers and prices. Underline key words in the questions.
\item \textbf{First listening:} do not write full sentences. Just catch the \cle{key words} (names, numbers, places, actions) and take short notes. Use abbreviations: \emph{MB, GB, FCFA, WhatsApp, USB}.
\item \textbf{Second listening:} check your answers, complete your notes, and verify the details you missed. Focus on the specific information required by each question.
\end{enumerate}
\end{methode}

\textbf{Listening text 1 (teacher reads aloud): Buying an Internet bundle.}\par
\emph{``Hello, I want to buy an Internet bundle for my smartphone. — Welcome! We have a daily bundle at 200 FCFA for 100 MB, and a monthly bundle at 5 000 FCFA for 2 GB. — I chat on WhatsApp and I browse websites every day, so I will take the monthly bundle, please. — Good choice. Dial star-one-five-zero and follow the instructions.''}

\textbf{Listening text 2: Sending files by WhatsApp.}\par
\emph{``Manka: Did you receive the photos of the workshop? — Ali: No, not yet. — Manka: I uploaded them on WhatsApp ten minutes ago. Check your connection. — Ali: Ah yes! I am downloading them now. The files are heavy, almost 15 MB. — Manka: Next time, I will copy them on your USB flash drive. It is faster!''}

\begin{exoresolu}[Listening comprehension — check your strategy]
\textbf{Questions (read before listening):} 1) How much does the monthly bundle cost? 2) Which two activities does the customer do every day? 3) Why does Manka prefer a USB flash drive next time?
\tcblower
\textbf{Solution:} 1) The monthly bundle costs \textbf{5 000 FCFA} for 2 GB (key word: a number + FCFA). 2) The customer \textbf{chats on WhatsApp} and \textbf{browses websites} every day (key words: chat, browse). 3) Because the USB flash drive is \textbf{faster} than downloading heavy files (15 MB) on WhatsApp. Notice how reading the questions first told us exactly which details to catch.
\end{exoresolu}

\courssub{Grammar Focus: \emph{used for + verb-ing}, \emph{used in + verb-ing} and \emph{used with + noun}}

Look at these two sentences from the listening texts:

\begin{itemize}
\item \emph{A USB flash drive \textbf{is used for storing} files.}
\item \emph{WhatsApp \textbf{is used with} a smartphone.}
\end{itemize}

Both sentences describe the \textbf{function} or the \textbf{use} of an object. But they answer two different questions:

\begin{propriete}[Expressing function, field of use and compatibility]
\begin{itemize}
\item \cle{used for + verb-ing} answers the question \emph{``What is it used \textbf{for}?''} — it expresses the \textbf{function} of the object.\\ Structure: \textbf{subject + be + used for + V-ing}.\\ Example: \emph{An MP3 player \textbf{is used for listening to} music.}
\item \cle{used in + verb-ing} expresses the \textbf{field or activity} where the object is used.\\ Structure: \textbf{subject + be + used in + V-ing}.\\ Example: \emph{Smartphones \textbf{are used in} teaching and learning.}
\item \cle{used with + noun} answers the question \emph{``What is it used \textbf{with}?''} — it expresses \textbf{compatibility}, the tool that goes together with it.\\ Structure: \textbf{subject + be + used with + noun}.\\ Example: \emph{This application \textbf{is used with} a smartphone.}
\end{itemize}
\end{propriete}

\begin{exemplebox}[A concrete example with numbers]
\emph{A 32 GB pen drive is used for storing about 8 000 songs in MP3 format.}\\
Here, \emph{is used for storing} gives the function (storing), and the number 8 000 makes the information precise and realistic. When you describe a gadget at the Probatoire or the Bac, always give its function with \emph{used for + V-ing}.
\end{exemplebox}

\begin{attention}[Do not confuse three different structures!]
The most frequent error is to mix these structures:
\begin{itemize}
\item \emph{I \textbf{used to play} football.} = past habit, finished now (\textbf{used to + infinitive}).
\item \emph{A pen drive \textbf{is used for storing} files.} = function (\textbf{be used for + V-ing}).
\item \emph{I \textbf{am used to using} a smartphone.} = I am accustomed to it (\textbf{be used to + V-ing}).
\end{itemize}
At the exam, ask yourself: \emph{function, past habit, or habit now?} Then choose the correct structure.
\end{attention}

\courssub{Classifying Activity: Match Each Gadget to Its Function}

\begin{exercice}[Match and build sentences]
Match each gadget (1--6) to its function (a--f), then write one full sentence with \emph{used for + V-ing} and one with \emph{used with + noun} for two gadgets of your choice.
\begin{enumerate}
\item USB flash drive \hspace{1cm} a. making free voice calls
\item MP3 player \hspace{1.55cm} b. sending formal messages
\item Viber \hspace{2.55cm} c. storing and carrying documents
\item e-mail \hspace{2.4cm} d. listening to music
\item smartphone \hspace{1.85cm} e. browsing websites
\item WhatsApp \hspace{1.95cm} f. chatting and sharing photos
\end{enumerate}
\end{exercice}

\begin{corrige}[Classifying activity]
Matching: 1--c, 2--d, 3--a, 4--b, 5--e, 6--f.\\
Example sentences: \emph{A USB flash drive is used for storing and carrying documents. It is used with a computer.} \emph{WhatsApp is used for chatting and sharing photos. It is used with a smartphone.}
\end{corrige}

\courssub{Statistics: Messaging Applications among Young Cameroonians}

Numbers help us speak about ICTs with precision. The bar chart below shows a realistic estimate of the use of messaging and social media applications by young Cameroonians.

\begin{popfigure}
\begin{center}
\begin{tikzpicture}
\begin{axis}[
ybar, bar width=0.9cm,
width=0.85\linewidth, height=6cm,
ymin=0, ymax=100,
ylabel={Users (\%)},
symbolic x coords={WhatsApp,Facebook,Viber},
xtick=data,
nodes near coords,
nodes near coords align={vertical},
every node near coord/.append style={font=\small\bfseries},
ymajorgrids=true, grid style={popline!40},
]
\addplot[fill=popblue, draw=popdark] coordinates {(WhatsApp,85) (Facebook,60) (Viber,15)};
\end{axis}
\end{tikzpicture}
\end{center}
\legende{Estimated use of messaging applications among young Cameroonians (fictitious but realistic data): WhatsApp 85\,\%, Facebook 60\,\%, Viber 15\,\%. WhatsApp is clearly the dominant application.}
\end{popfigure}

Read the chart in full sentences: \emph{``WhatsApp is used by about 85\,\% of young Cameroonians, while Viber is used by only 15\,\%.''} Notice how the passive voice, learned in the previous section, comes back naturally here.

\courssub{Speaking Practice: Mini-Dialogue about ICTs}

In pairs, practise this dialogue. Then replace the underlined parts with other gadgets or applications from the vocabulary list.

\begin{center}
\begin{tabularx}{\linewidth}{|X|}
\hline
\rowcolor{popblueL} \textbf{Dialogue: Choosing a gadget} \\
\hline
\textbf{A:} What is that \underline{smartphone} \textbf{used for}? \\
\textbf{B:} It \underline{is used for browsing} the Internet and \underline{chatting} on WhatsApp. \\
\textbf{A:} And what is it \underline{used with}? \\
\textbf{B:} It \underline{is used with} a \underline{4G SIM card} and a \underline{charger}. \\
\textbf{A:} Is it \underline{used in} education? \\
\textbf{B:} Yes, smartphones \underline{are used in} learning and research. \\
\hline
\end{tabularx}
\end{center}

\courssub{Discussion: Responsible Use of Social Media}

ICTs are powerful tools, but they also carry risks. In pairs, discuss these questions in English, using the vocabulary of the lesson:

\begin{itemize}
\item \textbf{Fake news:} \emph{Before you share a photo or a message on WhatsApp, what should you check?} Possible answer: \emph{``Before sharing, I check the source and I look for the information on a serious website.''}
\item \textbf{Cybercrime:} \emph{What dangers exist on the Internet?} Possible answer: \emph{``Some people steal passwords and money. We must never share our passwords or our Mobile Money codes.''}
\item \textbf{Time management:} \emph{How many hours per day do you chat or browse? Is it good for your studies?}
\item \textbf{Data protection:} \emph{Why should you not upload personal photos on public websites?}
\end{itemize}

\begin{aretenir}[What you must remember]
\begin{itemize}
\item ICTs are technologies used to create, store, exchange and use information in digital form.
\item Key vocabulary: www, website, MP3 player, USB flash/pen drive, WhatsApp, Viber, e-mail, smartphone, download, upload, browse, chat.
\item Listening strategy: read the questions first, catch key words, verify at the second listening.
\item \textbf{used for + V-ing} = function (\emph{What is it used for?}); \textbf{used in + V-ing} = field of use (\emph{Where is it used?}); \textbf{used with + noun} = compatibility (\emph{What is it used with?}).
\item Never confuse \emph{used to + infinitive} (past habit) with \emph{be used for + V-ing} (function) or \emph{be used to + V-ing} (being accustomed).
\end{itemize}
\end{aretenir}

\begin{exercice}[Consolidation — complete the sentences]
Complete with \emph{used for}, \emph{used in}, \emph{used with} or \emph{used to}:
\begin{enumerate}
\item A smartphone is \ldots\ browsing the Internet and chatting.
\item This application is \ldots\ a computer or a tablet.
\item When I was a child, I \ldots\ listen to the radio every evening.
\item An MP3 player is \ldots\ playing music.
\item E-mail is \ldots\ sending formal messages to teachers and employers.
\item Computers are \ldots\ education and research in modern schools.
\end{enumerate}
\end{exercice}

\begin{corrige}[Consolidation]
1) used for (function) ; 2) used with (compatibility) ; 3) used to (past habit, followed by the infinitive \emph{listen}) ; 4) used for (function) ; 5) used for (function) ; 6) used in (field of use).
\end{corrige}

You can now understand spoken texts about ICTs, name the main gadgets and applications, and describe their functions with precision. In the next section, \emph{Reading: Visiting a Workshop}, you will discover how these technologies are used and maintained in a real professional environment.

\coursec{Reading: Visiting a Workshop}

In the previous section, you listened to texts about the new Information and Communication Technologies: smartphones, USB flash drives, WhatsApp, Viber. But what happens when these gadgets break down? They are taken to a \cle{workshop}. In this section, you will read a text about a visit to a phone and computer repair workshop in Douala, learn the vocabulary of tools and machines, and practise expressing your \cle{preferences}.

\courssub{Key Vocabulary: The Workshop}

\begin{definition}[Workshop]
A \cle{workshop} (\emph{un atelier}) is a room or building where machines are used, repaired or maintained. In a technical school, each speciality has its own workshop: a metalwork workshop, an electrical workshop, a computer workshop.
\end{definition}

\begin{center}
\begin{tabularx}{\linewidth}{|l|X|X|}
\hline
\rowcolor{popblueL} \textbf{Word} & \textbf{French gloss} & \textbf{Example in context} \\
\hline
workshop & atelier & The students visited the repair workshop. \\
\hline
machine & machine & The machine cuts the metal sheets. \\
\hline
gadget & gadget, appareil & A smartphone is a useful gadget. \\
\hline
tool & outil & The technician uses small screwdrivers and other tools. \\
\hline
spare part & pièce de rechange & We ordered new spare parts for the printer. \\
\hline
maintenance & entretien, maintenance & Regular maintenance prevents breakdowns. \\
\hline
repair & réparation ; réparer & The repair of the screen costs \num{15000} FCFA. \\
\hline
to handle & manipuler, manier & Handle the motherboard with care. \\
\hline
to operate & faire fonctionner & Only trained students may operate this machine. \\
\hline
safety rules & règles de sécurité & Always follow the safety rules in the workshop. \\
\hline
\end{tabularx}
\end{center}

\begin{attention}[Handle with care!]
Do not confuse \cle{to handle} (to touch and use something carefully, \emph{manipuler}) and \cle{to operate} (to make a machine work, \emph{faire fonctionner}). You \emph{handle} a fragile screen; you \emph{operate} a drilling machine.
\end{attention}

\courssub{How to Read Efficiently: Skimming and Scanning}

\begin{methode}[Two steps to understand a text]
\begin{enumerate}
\item \textbf{Skimming} (\emph{lecture rapide}) : read the title, the first and the last sentences of each paragraph quickly. Goal: find the \cle{general idea}. Do not stop on unknown words.
\item \textbf{Scanning} (\emph{lecture ciblée}) : read again, looking for \cle{precise details}: names, dates, numbers, places. Use the questions (who? where? when?) to guide your eyes.
\end{enumerate}
\end{methode}

\courssub{Reading Text: A Visit to a Repair Workshop in Douala}

\textbf{Read the text first with skimming. What is it about?}

\begin{exemplebox}[Text: ``Our Visit to Fix-It Workshop'']
Last Tuesday, the students of Première Technique of Lycée Technique de Douala visited ``Fix-It Workshop'', a phone and computer repair workshop near Deïdo market. The visit started at 9 a.m. and finished at noon.

At the reception, the manager, Mr Etoundi, welcomed the group and presented the safety rules: do not touch the machines without permission, wear protective glasses near the soldering station, and keep the floor clear of cables.

Then the students moved to the repair bench, where four technicians were working. In this workshop, 15 phones are repaired every day. Screens are replaced more often than batteries, because many customers drop their phones. One technician showed the students how to handle a motherboard with antistatic gloves. Another one explained how he uses a multimeter to test the circuits.

Next, the group entered the storage room. Hundreds of spare parts were arranged on shelves: screens, batteries, chargers, keyboards and USB connectors. ``Good maintenance and original spare parts are the secret of a long-lasting repair,'' said Mr Etoundi.

Before leaving, the students asked questions. Many of them said they preferred practical work in the workshop to theory in the classroom. The visit ended with a group photo in front of the workshop.
\end{exemplebox}

\begin{popfigure}
\begin{tikzpicture}[scale=1, every node/.style={font=\small}]
% contour de l'atelier
\draw[thick, popdark] (0,0) rectangle (12,7);
% reception
\draw[fill=popblueL, draw=popblue] (0,4.5) rectangle (4,7);
\node at (2,5.9) {\textbf{1. Reception}};
\node[font=\footnotesize] at (2,5.3) {welcome desk};
% repair bench
\draw[fill=poporangeL, draw=poporange] (4,4.5) rectangle (12,7);
\node at (8,6.4) {\textbf{2. Repair bench}};
\node[font=\footnotesize] at (8,5.7) {4 technicians -- soldering station};
% petits bancs
\draw[fill=white, draw=poporange] (4.6,4.8) rectangle (5.8,5.4);
\draw[fill=white, draw=poporange] (6.4,4.8) rectangle (7.6,5.4);
\draw[fill=white, draw=poporange] (8.2,4.8) rectangle (9.4,5.4);
\draw[fill=white, draw=poporange] (10,4.8) rectangle (11.2,5.4);
% storage
\draw[fill=popgreenL, draw=popgreen] (0,0) rectangle (6,4.5);
\node at (3,3.9) {\textbf{3. Storage room}};
\node[font=\footnotesize] at (3,3.3) {spare parts on shelves};
\foreach \y in {0.6,1.4,2.2} {\draw[popgreen] (0.5,\y) -- (5.5,\y);}
% safety equipment
\draw[fill=poppinkL, draw=poppink] (6,0) rectangle (12,4.5);
\node at (9,3.9) {\textbf{4. Safety equipment}};
\node[font=\footnotesize] at (9,3.3) {fire extinguisher, gloves, glasses};
\draw[fill=white, draw=poppink] (7,1) circle (0.5);
\node[font=\footnotesize] at (7,1) {ext.};
\draw[fill=white, draw=poppink] (10.5,1) rectangle (11.5,2);
\node[font=\footnotesize] at (11,1.5) {kit};
% porte
\draw[thick, popdark] (5.2,0) arc[start angle=0, end angle=90, radius=0.9];
\draw[thick, popdark] (5.2,0) -- (6.1,0);
\node[font=\footnotesize, popdark] at (5.6,-0.4) {entrance};
% nord
\draw[->, thick, popdark] (12.6,6) -- (12.6,7);
\node[font=\footnotesize, popdark] at (12.6,7.2) {N};
\end{tikzpicture}
\legende{Sketch map of ``Fix-It Workshop'' (not to scale). 1: Reception. 2: Repair bench with four work stations. 3: Storage room for spare parts. 4: Safety equipment corner. The entrance is at the bottom.}
\end{popfigure}

\courssub{Comprehension: Finding the Information}

\begin{methode}[The WH-grid]
To check your understanding, complete the grid by scanning the text:
\begin{itemize}
\item \textbf{Who?} -- the people in the text.
\item \textbf{Where?} -- the places.
\item \textbf{When?} -- the time, the date.
\item \textbf{What?} -- the actions and events.
\item \textbf{How?} -- the way things are done.
\end{itemize}
\end{methode}

\begin{exoresolu}[Comprehension grid]
Complete the WH-grid for the text ``Our Visit to Fix-It Workshop''.
\tcblower
\begin{itemize}
\item \textbf{Who?} The students of Première Technique, their teachers, Mr Etoundi (the manager) and four technicians.
\item \textbf{Where?} At ``Fix-It Workshop'', near Deïdo market, in Douala.
\item \textbf{When?} Last Tuesday, from 9 a.m. to noon.
\item \textbf{What?} The students visited the reception, the repair bench and the storage room; they watched repairs and asked questions.
\item \textbf{How?} The technicians handle motherboards with antistatic gloves and test circuits with a multimeter.
\end{itemize}
\end{exoresolu}

\begin{exercice}[True or False?]
Read the text again and write \textbf{True} or \textbf{False}. Justify with a sentence from the text.
\begin{enumerate}
\item The workshop is in Yaoundé.
\item The visit lasted three hours.
\item Batteries are replaced more often than screens.
\item The students saw the storage room.
\item The students preferred theory to practical work.
\end{enumerate}
\end{exercice}

\begin{corrige}[Exercice : True or False?]
\begin{enumerate}
\item \textbf{False.} The workshop is near Deïdo market, in Douala.
\item \textbf{True.} It started at 9 a.m. and finished at noon (three hours).
\item \textbf{False.} Screens are replaced more often than batteries.
\item \textbf{True.} ``Next, the group entered the storage room.''
\item \textbf{False.} They preferred practical work in the workshop to theory in the classroom.
\end{enumerate}
\end{corrige}

\begin{exercice}[Find the equivalents]
Find in the text words or expressions that mean:
\begin{enumerate}
\item the person who manages the workshop (paragraph 2);
\item pieces used to replace broken parts (paragraph 4);
\item an instrument to measure electrical values (paragraph 3);
\item to stay for a long time without breaking (paragraph 4).
\end{enumerate}
\end{exercice}

\begin{corrige}[Exercice : Find the equivalents]
\begin{enumerate}
\item the manager (Mr Etoundi);
\item spare parts;
\item a multimeter;
\item long-lasting.
\end{enumerate}
\end{corrige}

\courssub{Grammar Focus: Expressing Preferences}

Imagine: your teacher offers you two tasks -- soldering a circuit or writing a report. You answer: ``I prefer soldering.'' You are expressing a \cle{preference}: a choice between two things.

\begin{propriete}[Structures of preference]
\cle{prefer} and \cle{would rather} express a choice:
\begin{itemize}
\item \textbf{prefer + noun} : I prefer practical work.
\item \textbf{prefer + V-ing} : I prefer repairing phones.
\item \textbf{prefer X to Y} : I prefer practical work \emph{to} theory. (X and Y must be parallel: two nouns or two V-ing.)
\item \textbf{would rather + base verb (without \emph{to})} : I would rather \emph{work} in the workshop than \emph{sit} in the classroom.
\item \textbf{like X better than Y} : I like laptops better than desktops.
\end{itemize}
\end{propriete}

\begin{exemplebox}[In the workshop]
\begin{itemize}
\item I prefer \textbf{using a laptop to} a desktop.
\item She prefers \textbf{repairing} screens. / She prefers \textbf{to repair} screens.
\item We would rather \textbf{learn} maintenance than \textbf{watch} a video.
\item They like the multimeter \textbf{better than} the oscilloscope.
\end{itemize}
\end{exemplebox}

\begin{attention}[Common error: ``prefer X to Y'']
Both ``I prefer \emph{to repair}'' and ``I prefer \emph{repairing}'' are correct. But with \textbf{prefer X to Y}, the two elements must be parallel: two nouns or two V-ing forms.
\begin{itemize}
\item Correct: I prefer \textbf{repairing} to \textbf{writing}. / I prefer \textbf{laptops} to \textbf{desktops}.
\item Wrong: I prefer repairing to \emph{write}. \quad Wrong: I prefer to repairing.
\item Wrong: I would rather \emph{to work}. $\rightarrow$ I would rather \textbf{work}.
\end{itemize}
\end{attention}

\begin{exoresolu}[Express your preference]
Complete with the correct form of the verb in brackets.
\begin{enumerate}
\item I prefer (use) \dots\ a laptop to a desktop.
\item The technician would rather (replace) \dots\ the battery than (buy) \dots\ a new phone.
\item My classmates prefer (practise) \dots\ in the workshop to (study) \dots\ theory.
\end{enumerate}
\tcblower
\begin{enumerate}
\item I prefer \textbf{using} a laptop to a desktop. (prefer V-ing to V-ing)
\item The technician would rather \textbf{replace} the battery than \textbf{buy} a new phone. (base verb, without \emph{to})
\item My classmates prefer \textbf{practising} in the workshop to \textbf{studying} theory. (two V-ing forms, parallel)
\end{enumerate}
\end{exoresolu}

\courssub{Tools and Their Functions: ``Used for \dots''}

Remember the structure from the listening section: \textbf{used for + V-ing} describes the function of a tool.

\begin{center}
\begin{tabularx}{\linewidth}{|l|X|X|}
\hline
\rowcolor{popblueL} \textbf{Tool} & \textbf{Function (used for\dots)} & \textbf{French gloss} \\
\hline
screwdriver & used for tightening and loosening screws & tournevis \\
\hline
soldering iron & used for joining electronic components & fer à souder \\
\hline
multimeter & used for measuring voltage and testing circuits & multimètre \\
\hline
pliers & used for holding and cutting wires & pince \\
\hline
antistatic gloves & used for handling sensitive components safely & gants antistatiques \\
\hline
\end{tabularx}
\end{center}

\courssub{Link with Your Speciality}

Every technical speciality has its own workshop and machines. Describe them using the vocabulary of this lesson and the preference structures.

\begin{exercice}[Your workshop, your preferences]
Answer in full sentences.
\begin{enumerate}
\item Name two machines or gadgets in the workshop of your speciality. What are they used for?
\item Which machine do you prefer operating? Why? (Use \emph{prefer} or \emph{would rather}.)
\item Write three sentences comparing two tools of your speciality, using \emph{prefer X to Y}, \emph{would rather} and \emph{like \dots\ better than}.
\end{enumerate}
\end{exercice}

\begin{corrige}[Exercice : Your workshop, your preferences (sample answers)]
\begin{enumerate}
\item In my workshop there is a drilling machine and a grinder. The drilling machine is used for making holes; the grinder is used for smoothing metal parts.
\item I prefer operating the drilling machine because it is precise and safe when you follow the safety rules.
\item I prefer the grinder to the file because it is faster. / I would rather use the multimeter than guess the voltage. / I like practical exercises better than long explanations.
\end{enumerate}
\end{corrige}

\begin{aretenir}[What to remember]
\begin{itemize}
\item A \cle{workshop} is a room or building where machines are used, repaired or maintained.
\item Read in two steps: \cle{skimming} for the general idea, \cle{scanning} for precise details (who, where, when, what, how).
\item Preferences: \textbf{prefer + noun/V-ing}, \textbf{prefer X to Y} (parallel forms), \textbf{would rather + base verb} (no \emph{to}), \textbf{like X better than Y}.
\item Functions of tools: \textbf{used for + V-ing}.
\end{itemize}
\end{aretenir}

In the next section, you will study punctuation and WH-/How questions, so that you can write a full report about your own workshop visit.

\coursec{Grammar Focus 2: Punctuation, WH- and How Questions}

In the previous section, we read a report about visiting a workshop in Douala. When the students interviewed the technician, they had to ask many questions: \emph{``What is this machine used for?''}, \emph{``How often do you maintain it?''}. They also had to write their report with correct punctuation. Asking good questions and punctuating correctly are two essential skills --- for a journalist at CRTV, for a technician writing a maintenance report, and for you at the Probatoire and Baccalauréat technique. This section gives you all the tools.

\courssub{English Punctuation: A Complete Review}

\begin{definition}[Punctuation]
\cle{Punctuation} (\emph{la ponctuation}) is the set of marks used in writing to separate sentences and their elements, and to clarify meaning. A wrongly placed mark can completely change the meaning of a sentence.
\end{definition}

Look at these two sentences:
\begin{itemize}
\item \emph{Let's eat, grandfather!} (We invite grandfather to eat.)
\item \emph{Let's eat grandfather!} (We want to eat grandfather!)
\end{itemize}
Only one comma, but a very different meaning! Punctuation saves lives --- and marks at the exam.

\begin{popfigure}
\begin{center}
\begin{tabularx}{\linewidth}{|c|l|X|}
\hline
\rowcolor{popblueL} \textbf{Mark} & \textbf{English name} & \textbf{Example} \\
\hline
. & full stop (\emph{point}) & The technician repairs phones. \\
\hline
, & comma (\emph{virgule}) & We saw lathes, drills, and presses. \\
\hline
; & semicolon (\emph{point-virgule}) & The workshop is noisy; however, it is well organised. \\
\hline
? & question mark (\emph{point d'interrogation}) & Where does he work? \\
\hline
! & exclamation mark (\emph{point d'exclamation}) & What a huge workshop! \\
\hline
' & apostrophe (\emph{apostrophe}) & It's the manager's office. \\
\hline
: & colon (\emph{deux-points}) & Bring these tools: a hammer, a screwdriver and pliers. \\
\hline
`` '' & quotation marks (\emph{guillemets}) & The boss said, ``Safety first.'' \\
\hline
A & capital letter (\emph{majuscule}) & On Monday, Ali speaks English. \\
\hline
\end{tabularx}
\end{center}
\legende{The main English punctuation marks, their names and one example for each.}
\end{popfigure}

\begin{aretenir}[Capital letters in English: the differences with French]
English uses capital letters (\emph{majuscules}) in more situations than French. Always write a capital letter for:
\begin{itemize}
\item the pronoun \cle{I}: \emph{My friend and I visited the workshop.} (never \emph{i})
\item \cle{nationalities} and \cle{languages}: \emph{English, French, Cameroonian, Nigerian} (\emph{anglais, camerounais} --- no capital in French, but obligatory in English!)
\item \cle{days} and \cle{months}: \emph{Monday, Tuesday, January, February} (but seasons stay small: \emph{rainy season})
\item the first word of a sentence and proper nouns: \emph{CRTV, Mount Cameroon, Yaoundé}.
\end{itemize}
\end{aretenir}

\begin{savaistu}[Punctuation spacing: English vs French]
In French, you put a space \textbf{before} double punctuation marks (\emph{? ! : ;}) and inside guillemets \emph{« ... »}. In English, \textbf{no space} before these marks, and quotation marks hug the text: \emph{``Safety first.''} not \emph{`` Safety first . ''}. This is a frequent error in exams.
\end{savaistu}

\begin{attention}[The apostrophe: two different jobs]
Do not confuse the two uses of the apostrophe:
\begin{enumerate}
\item \cle{Contraction}: \emph{it's} = it is; \emph{don't} = do not; \emph{I'm} = I am.
\item \cle{Possession}: \emph{the technician's tools} (the tools of the technician); \emph{students' reports} (the reports of the students --- apostrophe after the plural \emph{s}).
\end{enumerate}
Trap: \emph{its} (possessive adjective, no apostrophe) $\neq$ \emph{it's} (it is). \emph{The machine has lost its power. It's broken.}
\end{attention}

\courssub{WH-Questions: Asking for Information}

Imagine the situation: you are in the workshop and you want information. A yes/no question (\emph{Do you like your job?}) gives you only ``yes'' or ``no''. To get real information --- a name, a place, a time, a reason --- you need a WH-question.

\begin{definition}[WH-question]
A \cle{WH-question} (\emph{question partielle}) is a question introduced by an interrogative word (What, Where, When, Who, Why, Which, Whose, How) that \textbf{cannot be answered by yes or no}. It asks for a specific piece of information.
\end{definition}

\begin{center}
\begin{tabularx}{\linewidth}{|l|l|X|}
\hline
\rowcolor{popblueL} \textbf{Question word} & \textbf{Asks about...} & \textbf{Example} \\
\hline
What & a thing, an action & What is this machine used for? \\
\hline
Where & a place (\emph{lieu}) & Where do you work? \\
\hline
When & time (\emph{temps}) & When did you start this job? \\
\hline
Who & a person & Who repairs the computers? \\
\hline
Why & a reason (\emph{raison}) & Why do you wear gloves? \\
\hline
Which & a choice between things & Which tool is the most useful? \\
\hline
Whose & possession (\emph{possession}) & Whose helmet is this? \\
\hline
\end{tabularx}
\end{center}

\begin{methode}[Building a WH-question]
Follow these steps, in this order:
\begin{enumerate}
\item Choose the correct \textbf{question word} (What, Where, When...).
\item Add the \textbf{auxiliary}: \emph{do / does / did}, or \emph{be}, or a modal (\emph{can, must}...).
\item Add the \textbf{subject}.
\item Add the \textbf{base verb} (\emph{base verbale}, without -s, without -ed).
\end{enumerate}
Formula: \textbf{Question word + auxiliary + subject + base verb (+ complement)?}
\par Example: \emph{Where + does + he + work?} --- \emph{When + did + they + visit + the workshop?}
\end{methode}

\begin{attention}[The most frequent error: the ``s'' of the 3rd person]
After \emph{does}, the main verb takes the \textbf{base form}: the \emph{-s} disappears because \emph{does} already carries it.
\begin{itemize}
\item Correct: \emph{Where does he work\textbf{?}}
\item Wrong: \emph{Where does he work\textbf{s}?}
\end{itemize}
Same rule with \emph{did}: \emph{What did you see?} (not \emph{saw}). The past is already inside \emph{did}.
\end{attention}

\courssub{How Questions: Quantity, Frequency, Distance, Duration}

\cle{How} combines with adjectives or adverbs to ask precise questions:

\begin{center}
\begin{tabularx}{\linewidth}{|l|l|X|}
\hline
\rowcolor{popblueL} \textbf{Expression} & \textbf{Asks about...} & \textbf{Example} \\
\hline
How & manner, state (\emph{manière}) & How does this machine work? \\
\hline
How many & countable quantity (\emph{dénombrable}) & How many tools do you have? \\
\hline
How much & uncountable quantity (\emph{indénombrable}) & How much does it cost? \\
\hline
How often & frequency (\emph{fréquence}) & How often do you listen to the radio? \\
\hline
How long & duration (\emph{durée}) & How long does the repair take? \\
\hline
How far & distance & How far is the workshop from school? \\
\hline
How old & age & How old is this generator? \\
\hline
\end{tabularx}
\end{center}

\begin{attention}[How much vs How many]
This is a classic exam trap:
\begin{itemize}
\item \cle{How many} + \textbf{countable} noun (plural): \emph{How many phones, tools, channels, students...?}
\item \cle{How much} + \textbf{uncountable} noun: \emph{How much money, information, time, oil...?}
\end{itemize}
\emph{Information} and \emph{money} are uncountable in English: never \emph{informations}, never \emph{how many money}.
\end{attention}

\begin{exemplebox}[A journalist's questions]
\begin{itemize}
\item \emph{How many channels does CRTV broadcast?} --- It broadcasts several national and regional programmes.
\item \emph{How often do you listen to the radio?} --- Every morning before work.
\item \emph{How long have you worked at CRTV?} --- For ten years.
\end{itemize}
\end{exemplebox}

\begin{popfigure}
\begin{center}
\begin{tikzpicture}[
grow=right,
level distance=3.4cm,
level 1/.style={sibling distance=2.8cm},
level 2/.style={sibling distance=1.6cm},
every node/.style={font=\small},
qnode/.style={rectangle, rounded corners, draw=popblue, fill=popblueL, align=center, inner sep=3pt},
anode/.style={rectangle, rounded corners, draw=popgreen, fill=popgreenL, align=center, inner sep=3pt},
edge from parent/.style={draw=popmuted, thick}]
\node[qnode, align=center]{What do you\\want to know?}
  child { node[qnode]{quantity?} 
    child { node[anode, align=center]{uncountable noun\\$\to$ \textbf{how much}} }
    child { node[anode, align=center]{countable noun\\$\to$ \textbf{how many}} }
  }
  child { node[anode]{reason $\to$ \textbf{why}} }
  child { node[anode]{place $\to$ \textbf{where}} }
  child { node[anode]{person $\to$ \textbf{who}} }
  child { node[anode]{time $\to$ \textbf{when}} }
  child { node[anode]{choice $\to$ \textbf{which}} }
  child { node[anode]{possession $\to$ \textbf{whose}} };
\end{tikzpicture}
\end{center}
\legende{Decision tree: choosing the correct question word. First identify what you want to know (person, place, reason, time, choice, possession, quantity), then apply the rule --- for quantities, check if the noun is countable or uncountable.}
\end{popfigure}

\courssub{Subject Questions vs Object Questions}

Compare these two sentences carefully:

\begin{itemize}
\item \emph{Who repaired the phone?} --- \textbf{The technician} repaired it. (We ask about the \cle{subject}.)
\item \emph{Who did you call?} --- I called \textbf{the technician}. (We ask about the \cle{object} --- \emph{le complément}.)
\end{itemize}

\begin{propriete}[Two structures for WH-questions]
\begin{enumerate}
\item \textbf{Question about the subject}: Who/What replaces the subject. \textbf{No auxiliary, no inversion}; the verb keeps its normal form (with -s in the present).
\par \emph{Who repairs the machines? --- The engineer does.}
\par \emph{What happened yesterday? --- A machine broke down.}
\item \textbf{Question about the object or another element}: use the normal structure with auxiliary.
\par \emph{Who did you call? --- I called the engineer.}
\par \emph{What does this machine produce? --- It produces metal parts.}
\end{enumerate}
\end{propriete}

\begin{attention}[No ``do'' in a subject question]
Wrong: \emph{Who does repair the phone?} Correct: \emph{Who repairs the phone?} When the question word \emph{is} the subject, the auxiliary \emph{do/does/did} is not used. Test: if the answer replaces the question word as subject (\emph{The technician repaired it}), then it is a subject question.
\end{attention}

\begin{exoresolu}[Building questions for an interview]
Énoncé --- A student interviews a workshop technician. Write the question corresponding to each answer, using the underlined element.
\begin{enumerate}
\item \emph{The workshop opens at 7.30 a.m.} (When...)
\item \emph{Mr Etoundi repairs the generators.} (Who... --- question about the subject)
\item \emph{She uses a screwdriver to open the phone.} (What... --- question about the object)
\item \emph{The workshop has fifteen machines.} (How many...)
\end{enumerate}
\tcblower
Solution détaillée :
\begin{enumerate}
\item The underlined element is a time $\to$ \emph{When}. Present simple, subject ``the workshop'' (3rd person) $\to$ \emph{does} + base verb. \textbf{When does the workshop open?}
\item ``Mr Etoundi'' is the subject of ``repairs'' $\to$ subject question: no auxiliary, verb keeps the -s. \textbf{Who repairs the generators?}
\item ``a screwdriver'' is the object of ``uses'' $\to$ auxiliary \emph{does}, base verb \emph{use}. \textbf{What does she use to open the phone?}
\item ``fifteen'' is a countable quantity (machines) $\to$ \emph{How many} + plural noun + \emph{does} + subject + base verb. \textbf{How many machines does the workshop have?}
\end{enumerate}
\end{exoresolu}

\courssub{Application: Interviewing a Professional}

Here is a model set of ten interview questions you could ask a workshop technician or a CRTV journalist. Notice the variety of question words and the correct structures.

\begin{exemplebox}[Ten interview questions]
\begin{enumerate}
\item What is your job exactly?
\item Where do you work?
\item When did you start this career?
\item Why did you choose this profession?
\item How many people work in your team?
\item How much time do you spend on one repair?
\item How often do you maintain the machines?
\item How far is your workshop from your home?
\item Who trained you for this job? (subject question)
\item Which tool do you prefer, and why?
\end{enumerate}
\end{exemplebox}

\begin{exercice}[Punctuation and capital letters]
Rewrite these sentences with correct punctuation and capital letters.
\begin{enumerate}
\item the technician said safety first
\item i visited a workshop in douala on monday
\item what a modern workshop
\item bring these tools a hammer a drill and pliers
\item its the managers office
\end{enumerate}
\end{exercice}

\begin{corrige}[Exercice : Punctuation and capital letters]
\begin{enumerate}
\item The technician said, ``Safety first.''
\item I visited a workshop in Douala on Monday.
\item What a modern workshop!
\item Bring these tools: a hammer, a drill and pliers.
\item It's the manager's office.
\end{enumerate}
\end{corrige}

\begin{exercice}[Build the WH-question]
Write the question that produced each answer. The underlined words give the answer.
\begin{enumerate}
\item \emph{The radio studio is in Yaoundé.} (Where...)
\item \emph{Mrs Ngo Bell presents the news.} (Who... --- subject question)
\item \emph{He checks the machines twice a week.} (How often...)
\item \emph{The repair costs 5\,000 francs.} (How much...)
\item \emph{They visited the workshop last Friday.} (When...)
\end{enumerate}
\end{exercice}

\begin{corrige}[Exercice : Build the WH-question]
\begin{enumerate}
\item Where is the radio studio?
\item Who presents the news? (subject question: no auxiliary, verb with -s)
\item How often does he check the machines?
\item How much does the repair cost?
\item When did they visit the workshop?
\end{enumerate}
\end{corrige}

\begin{exercice}[Subject or object question?]
Decide if the question asks about the subject or the object, then write the correct question for each answer.
\begin{enumerate}
\item \emph{The engineer designed the new circuit.} (Who...?)
\item \emph{The technician called the engineer.} (Who...?)
\item \emph{A power cut stopped the production.} (What...?)
\item \emph{The manager checked the safety rules.} (What...?)
\end{enumerate}
\end{exercice}

\begin{corrige}[Exercice : Subject or object question?]
\begin{enumerate}
\item Who designed the new circuit? (subject question: \emph{The engineer} is subject)
\item Who did the technician call? (object question: \emph{the engineer} is object)
\item What stopped the production? (subject question: \emph{A power cut} is subject)
\item What did the manager check? (object question: \emph{the safety rules} is object)
\end{enumerate}
\end{corrige}

\begin{exercice}[Your turn: prepare an interview]
Write ten questions to interview a CRTV journalist or a workshop technician in your town. Use at least: two \emph{How many/much} questions, one \emph{How often}, one \emph{How long}, one subject question with \emph{Who}, and one \emph{Why} question. Check every auxiliary and every base verb.
\end{exercice}

\begin{corrige}[Exercice : Your turn --- elements of correction]
Accept any grammatically correct set. Example of a correct production:
\begin{enumerate}
\item What does a journalist do every day?
\item Where do you find your information?
\item When did you join CRTV?
\item Why do you like this job?
\item Who writes the news bulletins? (subject question)
\item How many reports do you prepare per week?
\item How much time do you spend in the field?
\item How often do you broadcast live?
\item How long does an interview usually last?
\item Which programme do you prefer?
\end{enumerate}
Check: base verb after \emph{does/did}; no auxiliary in the subject question; \emph{how much} with uncountable nouns only.
\end{corrige}

\begin{aretenir}[Summary: what you must remember]
\begin{itemize}
\item Capital letters: \emph{I}, nationalities, languages, days, months, sentence beginnings.
\item WH-question = question word + auxiliary + subject + \textbf{base verb}; it cannot be answered by yes/no.
\item After \emph{does} and \emph{did}, the verb loses its -s and its -ed.
\item \emph{How many} = countable; \emph{How much} = uncountable.
\item Subject question (\emph{Who repairs...?}): no auxiliary, verb agrees with the unknown subject.
\item Punctuation spacing: no space before ? ! : ; in English (unlike French).
\end{itemize}
You are now ready for the final section: writing your own report about a workshop visit, with correctly punctuated sentences and well-built questions.
\end{aretenir}

\coursec{Writing \& Integration: Reporting a Workshop Visit}

In the previous sections, you learnt how to use \cle{punctuation} correctly and how to ask \cle{WH- and How questions}. Now it is time to use \emph{everything} you have studied in this lesson --- the passive voice, \emph{used for + verb-ing}, preferences, media and workshop vocabulary, punctuation and questions --- to do what real journalists and technicians do: \textbf{write a report}. In this final section, you will learn how to write a good composition about a workshop visit, how to check and correct your work, and how to succeed in the integration activities and the final evaluation.

\courssub{The Structure of a Composition}

A good composition is like a good journey: it has a clear \textbf{beginning}, a rich \textbf{middle} and a satisfying \textbf{end}. Examiners at the Probatoire technique always look for these three parts.

\begin{definition}[The three parts of a composition]
A \cle{composition} is a short organised text. It has three parts:
\begin{itemize}
\item \textbf{Introduction}: answers the questions \emph{Where? When? Who?} (e.g. \emph{Last Friday, our class visited a phone repair workshop in Douala.})
\item \textbf{Body}: describes the \emph{machines seen}, the \emph{activities observed} and your \emph{preferences} (e.g. \emph{I preferred the soldering station because\ldots})
\item \textbf{Conclusion}: states the \emph{lesson learnt} (e.g. \emph{This visit taught me that machines must be maintained regularly.})
\end{itemize}
\end{definition}

\begin{popfigure}
\begin{center}
\begin{tikzpicture}[scale=0.9]
% Introduction layer (top)
\fill[popblue!80] (0,0) -- (10,0) -- (9,-1.5) -- (1,-1.5) -- cycle;
\node[white,font=\bfseries\large] at (5,-0.75) {INTRODUCTION: Where? When? Who?};
% Body layer (middle)
\fill[poporange!85] (1.5,-2) -- (8.5,-2) -- (7.8,-4) -- (2.2,-4) -- cycle;
\node[white,font=\bfseries\large] at (5,-3) {BODY: machines, activities, preferences};
% Conclusion layer (bottom)
\fill[popgreen!80] (3,-4.5) -- (7,-4.5) -- (6.5,-5.8) -- (3.5,-5.8) -- cycle;
\node[white,font=\bfseries\large] at (5,-5.15) {CONCLUSION};
\node[popdark,font=\small] at (5,-6.2) {the lesson learnt};
% Arrows on sides
\draw[popdark, thick, ->] (0.5,-1.5) -- (0.5,-2);
\draw[popdark, thick, ->] (9.5,-1.5) -- (9.5,-2);
\draw[popdark, thick, ->] (2,-4) -- (2,-4.5);
\draw[popdark, thick, ->] (8,-4) -- (8,-4.5);
\end{tikzpicture}
\end{center}
\legende{The funnel structure of a composition: from general information (introduction) to specific details (body) and finally to one clear idea (conclusion).}
\end{popfigure}

\courssub{Using What You Have Learnt}

Your composition will be excellent if you \emph{recycle} the grammar and vocabulary of this lesson:

\begin{itemize}
\item \textbf{Passive voice} to describe processes: \emph{The phones are opened carefully. The screens are tested. Broken parts are replaced.}
\item \textbf{Used for + verb-ing} to describe functions: \emph{A multimeter is used for measuring voltage. A USB flash drive is used for storing data.}
\item \textbf{Used with + noun} to describe materials/tools: \emph{A soldering iron is used with tin. A computer is used with diagnostic software.}
\item \textbf{Preferences}: \emph{I prefer the radio studio to the repair bench. I like soldering more than cleaning.}
\item \textbf{Linking words} (\cle{connectors}) to organise your ideas: \emph{first, then, after that, finally, moreover, however, in conclusion}.
\item \textbf{Correct punctuation}: capital letter at the beginning of every sentence, full stop at the end, commas between items in a list, question marks for questions.
\end{itemize}

\begin{methode}[How to write a composition: plan, draft, check, copy]
\begin{enumerate}
\item \textbf{Plan (brainstorming)}: write down your ideas in note form --- place, date, people, machines, activities, preferences, lesson learnt.
\item \textbf{Draft}: write your first version in three paragraphs (introduction, body, conclusion). Use linking words.
\item \textbf{Check}: read your draft again. Verify the \emph{grammar} (passive voice, tenses, \emph{used for/with}), the \emph{punctuation} (capital letters, full stops, commas) and the \emph{linking words}.
\item \textbf{Final copy}: rewrite the corrected composition neatly.
\end{enumerate}
\end{methode}

\begin{attention}[Frequent errors to avoid]
Many candidates lose easy marks because of:
\begin{itemize}
\item sentences \textbf{without a final full stop};
\item sentences \textbf{without a capital letter} at the beginning;
\item forgetting the \textbf{-ed} of the past participle in the passive: \emph{the screen is replace} instead of \emph{the screen is \textbf{replaced}};
\item mixing up \emph{used for + verb-ing} and \emph{used with + noun}: \emph{A soldering iron is used for \textbf{soldering}} (function) but \emph{it is used with \textbf{tin}} (material);
\item writing \emph{used for measure} instead of \emph{used for \textbf{measuring}};
\item forgetting the auxiliary \emph{be} in the passive: \emph{the phones repaired} instead of \emph{the phones \textbf{are} repaired}.
\end{itemize}
\end{attention}

\courssub{A Guided Writing Task}

\begin{exemplebox}[A typical exam topic]
\emph{Write a 150-word composition about your visit to a phone repair workshop, saying what machines you saw and what they are used for.}
\par\medskip
\textbf{Plan (step 1):} Introduction: last Tuesday / class of Première technique / Mballa workshop, Yaoundé. Body: soldering iron (used for joining wires), multimeter (used for measuring voltage), computer (used for diagnosing faults); I preferred the multimeter. Conclusion: machines must be handled carefully.
\end{exemplebox}

\begin{exoresolu}[Model composition (about 150 words)]
Write the composition from the plan above.
\tcblower
\textbf{Our Visit to a Phone Repair Workshop}\par
Last Tuesday, the students of Première technique visited the Mballa phone repair workshop in Yaoundé. We were welcomed by the chief technician, Mr Ngo Bassa.\par
First, we were shown the soldering bench. A soldering iron is used for joining wires and small components. Then, we observed a multimeter: it is used for measuring voltage and testing batteries. After that, damaged phones were opened and their screens were tested on a computer, which is used for diagnosing faults. Moreover, spare parts are stored carefully in labelled boxes. I preferred the multimeter because it is precise and easy to handle; however, some students liked the soldering iron more.\par
In conclusion, this visit taught us that machines and gadgets must be handled and maintained carefully. It was a very useful experience for our future jobs.\par
(\emph{about 150 words; passive voice, ``used for'', preferences, connectors and punctuation are all used.})
\end{exoresolu}

\courssub{Peer Correction: The Evaluation Grid}

After drafting, exchange your composition with a classmate and correct it with the grid below. This is called \cle{peer correction}.

\begin{center}
\begin{tabularx}{\linewidth}{|l|X|c|}
\hline
\rowcolor{popblueL} \textbf{Criterion} & \textbf{What to check} & \textbf{Mark} \\
\hline
Grammar & passive voice correct, tenses, \emph{used for + verb-ing}, \emph{used with + noun}, agreement & /10 \\
\hline
Vocabulary & media and workshop words used correctly and varied & /5 \\
\hline
Punctuation & capital letters, full stops, commas, question marks & /5 \\
\hline
\textbf{Total} & & \textbf{/20} \\
\hline
\end{tabularx}
\end{center}

\begin{exercice}[Peer correction practice]
Your partner wrote: \emph{``last friday we visit a workshop the phones is repaired by the technician a multimeter is used for measure voltage''}. Find and correct \textbf{six} errors (punctuation, tense, passive, \emph{used for}).
\end{exercice}

\begin{corrige}[Peer correction practice]
\emph{``\textbf{L}ast \textbf{F}riday\textbf{, we \textbf{visited}} a workshop\textbf{.} The phones \textbf{were} repaired by the technician\textbf{.} A multimeter is used for \textbf{measuring} voltage\textbf{.}''}\par
Errors corrected: (1) capital letter \emph{Last}; (2) capital letter \emph{Friday}; (3) past tense \emph{visited}; (4) passive plural \emph{were repaired}; (5) \emph{used for measuring} (verb-ing); (6) missing full stops and comma added.
\end{corrige}

\courssub{Integration Activities}

\begin{definition}[Integration activity]
An \cle{integration activity} is a complex task that combines several skills (listening, reading, speaking, writing) in a real-life situation. It shows that you can \emph{use} English, not only \emph{know} it.
\end{definition}

\textbf{Complex task: ``The Young Reporter''.} Work in groups of four.
\begin{enumerate}
\item \textbf{Listening}: listen to a short text about a student who uses WhatsApp and a USB flash drive to prepare a report. Answer five WH-questions.
\item \textbf{Reading}: read a short article about a visit to a radio station workshop. Underline the passive sentences and the machines mentioned.
\item \textbf{Speaking}: interview the ``technician'' (a classmate) with WH- and How questions: \emph{What is this machine used for? How is it maintained?}
\item \textbf{Writing}: individually, write a 150-word composition reporting the visit, following the plan--draft--check--copy method.
\end{enumerate}

\courssub{Evaluation and Remediation}

\begin{exercice}[Summative test (end of lesson)]
\textbf{A. Grammar (/6):} Put the verbs into the passive voice: (1) The technician (repair) the radio every week. (2) The data (store) on a USB flash drive. (3) Complete: A printer is used for \ldots\ (print).\par
\textbf{B. Vocabulary (/4):} Match: \emph{broadcast, headline, soldering iron, pen drive} with their definitions.\par
\textbf{C. Reading (/4):} Read a short text about visiting a workshop and answer two WH-questions and two True/False questions.\par
\textbf{D. Writing (/6):} Write a 100-word paragraph about your favourite media or machine, using at least one passive sentence and one preference.
\end{exercice}

\begin{corrige}[Summative test --- extracts]
\textbf{A.} (1) The radio \emph{is repaired} by the technician every week. (2) The data \emph{is stored} on a USB flash drive. (3) A printer is used for \emph{printing}.\par
\textbf{B.} \emph{broadcast} = send out programmes; \emph{headline} = main title; \emph{soldering iron} = tool for joining metal; \emph{pen drive} = portable storage device.\par
\textbf{D.} Example: \emph{I prefer the radio to television. News is broadcast every hour, and songs are played between the programmes. I like listening while I work.}
\end{corrige}

\textbf{Remediation activities} (for students who need support):
\begin{itemize}
\item \textbf{Grammar support}: rewrite five active sentences in the passive voice with the teacher's help, then alone.
\item \textbf{Punctuation support}: copy a short paragraph and add all the missing capital letters and full stops.
\item \textbf{Vocabulary support}: label pictures of ten machines and gadgets, then write one sentence with \emph{used for} for each.
\item \textbf{Writing support}: complete a guided composition with gaps (linking words and verbs given in brackets).
\end{itemize}

\begin{savaistu}[Writing counts at the Probatoire technique!]
At the Probatoire and Baccalauréat technique in Cameroon, written expression represents an important part of the English paper. Punctuation and grammar are \emph{explicitly} marked: a composition with good ideas but no full stops and no capital letters can lose many points. Training with the plan--draft--check--copy method is the best way to succeed!
\end{savaistu}

\begin{aretenir}[Key points of the section]
\begin{itemize}
\item A composition has three parts: \textbf{introduction} (where/when/who), \textbf{body} (machines, activities, preferences), \textbf{conclusion} (lesson learnt).
\item Recycle the lesson's tools: \textbf{passive voice}, \emph{used for + verb-ing}, \emph{used with + noun}, preferences, linking words (\emph{first, then, after that, finally, moreover, however, in conclusion}).
\item Follow the method: \textbf{plan $\rightarrow$ draft $\rightarrow$ check $\rightarrow$ final copy}.
\item Every sentence begins with a \textbf{capital letter} and ends with a \textbf{full stop}.
\item An \textbf{integration activity} combines listening, reading, speaking and writing in a real-life task.
\end{itemize}
\end{aretenir}

\newpage

\coursec{Activité d'intégration}

\coursec{Media and Communication}

\courssub{Activity Objectives}

By the end of this integration activity, you will be able to:
\begin{itemize}
\item exchange information about the media and new technologies in an authentic situation (\cle{speaking});
\item prepare and ask \cle{WH-} and \cle{How questions} during an interview;
\item write a short radio report using the \cle{passive voice}, \emph{used for + verb-\emph{ing}} and \emph{used with + noun}, an expression of \cle{preference}, and correct \cle{punctuation};
\item present the report orally, as on the radio;
\item evaluate the work of other groups with a grid and correct your own mistakes (\cle{remediation}).
\end{itemize}

\begin{aretenir}[Competences integrated]
Speaking + Listening + Grammar (passive voice, \emph{used for}, preferences, WH-questions) + Writing. All the skills of the lesson are used in one real-life task: producing a radio report.
\end{aretenir}

\courssub{Situation}

\textbf{Radio Balafon needs you!}

Radio Balafon, a community radio station in Douala, broadcasts a weekly programme for young people called \emph{``Youth and Jobs''}. This week, the journalists want to show how technical schools prepare students for real jobs. Your class has been chosen: you are the \cle{reporters} (\emph{les reporters})!

Yesterday, your team visited the \textbf{maintenance workshop} (\emph{l'atelier de maintenance}) of your technical high school, where Mr. Essono, the workshop manager (\emph{le chef d'atelier}), repairs computers, phones and printers for the school and the neighbourhood.

Here is the data your team collected during the visit:

\begin{center}
\begin{tabularx}{\linewidth}{|l|X|}\hline
\rowcolor{popblueL} \textbf{Item} & \textbf{Data collected} \\ \hline
Machines in the workshop & 3 soldering stations, 2 diagnostic computers, 1 ultrasonic cleaner, 1 drilling machine \\ \hline
Devices repaired per week & 35 mobile phones, 12 laptops, 8 printers \\ \hline
Most frequent phone problems & broken screens (40 \%), battery problems (25 \%), charging ports (20 \%), software problems (15 \%) \\ \hline
Tools used & screwdrivers, multimeter, soldering iron, USB flash drives with repair software \\ \hline
Workers & 4 technicians and 6 students on internship (\emph{stagiaires}) \\ \hline
Working hours & Monday to Saturday, 8 a.m. to 5 p.m. \\ \hline
\end{tabularx}
\end{center}

Mr. Essono also told you: \emph{``The multimeter is used for testing electric circuits. The soldering iron is used with tin wire. I prefer repairing laptops because the work is cleaner and better paid.''}

Your mission: produce a \textbf{150-word radio report} about this visit, ready to be broadcast on Radio Balafon.

\courssub{Instructions}

Work in groups of four. Follow the steps in order.

\begin{enumerate}
\item \textbf{Prepare the interview.} Write \textbf{8 questions} for Mr. Essono: at least 5 \emph{WH-questions} (What, Where, When, Why, Which, Who) and at least 2 \emph{How questions} (How many, How often, How long\ldots). Use the data table to imagine what you still want to know.
\item \textbf{Conduct the interview.} One student plays Mr. Essono, the others are the journalists. Ask your questions, listen carefully and take notes. The ``manager'' must answer using the data of the situation.
\item \textbf{Write the radio report (about 150 words).} Your report must contain:
\begin{itemize}
\item a catchy introduction (who? where? when?);
\item at least \textbf{3 sentences in the passive voice};
\item at least \textbf{2 structures} with \emph{used for + verb-\emph{ing}} or \emph{used with + noun};
\item at least \textbf{1 expression of preference} (\emph{prefer, would rather, like \ldots\ better});
\item correct \textbf{punctuation}: capital letters, full stops, commas, question marks in quoted questions.
\end{itemize}
\item \textbf{Record or present your report} in front of the class, with a clear voice, like real radio journalists. One student introduces, one reads the report, one plays short interview extracts.
\item \textbf{Evaluate another group} with the grid below, then correct your own production after the feedback (\emph{remédiation}).
\end{enumerate}

\begin{methode}[Evaluation grid (10 points)]
\begin{center}
\begin{tabularx}{\linewidth}{|X|c|}\hline
\rowcolor{popgreenL} \textbf{Criterion} & \textbf{Points} \\ \hline
8 correct WH-/How questions prepared & 2 \\ \hline
At least 3 correct passive sentences & 2 \\ \hline
2 correct \emph{used for / used with} structures & 2 \\ \hline
1 expression of preference + correct punctuation & 2 \\ \hline
Clear oral presentation (voice, rhythm, pronunciation) & 2 \\ \hline
\end{tabularx}
\end{center}
\end{methode}

\courssub{Resources}

\begin{aretenir}[Grammar toolbox]
\begin{itemize}
\item \textbf{Passive voice}: \emph{subject + be + past participle (+ by + agent)}. Ex.: \emph{Thirty-five phones are repaired every week.}
\item \textbf{Used for + verb-\emph{ing}}: the purpose of a tool. Ex.: \emph{A multimeter is used for testing circuits.}
\item \textbf{Used with + noun}: what a tool works together with. Ex.: \emph{The soldering iron is used with tin wire.}
\item \textbf{Preferences}: \emph{I prefer laptops to phones. / I would rather repair laptops. / I like laptops better.}
\item \textbf{WH-/How questions}: \emph{question word + auxiliary + subject + verb?} Ex.: \emph{How many phones do you repair every week?}
\item \textbf{Punctuation}: capital letter at the beginning, full stop at the end, comma in lists, question mark after a direct question, quotation marks for reported speech.
\end{itemize}
\end{aretenir}

\begin{attention}[Common mistakes to avoid]
\begin{itemize}
\item Do not forget \emph{be} in the passive: \emph{phones are repaired}, not \emph{phones repaired}.
\item After \emph{used for}, always use \emph{verb-\emph{ing}}: \emph{used for testing}, not \emph{used for test}.
\item Word order in questions: \emph{How many phones \textbf{do you repair}?}, not \emph{How many phones you repair?}
\item Every direct question inside your report ends with a \textbf{question mark}.
\end{itemize}
\end{attention}

\courssub{Procedure and Model Answer}

\begin{exoresolu}[Model production: the Radio Balafon report]
\textbf{Step 1 -- The 8 interview questions.}\par
1. \emph{What machines do you use in this workshop?}\par
2. \emph{How many phones do you repair every week?}\par
3. \emph{Which problem is the most frequent?}\par
4. \emph{How long does a screen repair take?}\par
5. \emph{Why do you prefer repairing laptops?}\par
6. \emph{Who works with you in the workshop?}\par
7. \emph{When do you open the workshop?}\par
8. \emph{How often do students come here for internships?}\par
\textit{Comment: questions 1, 3, 5, 6, 7 are WH-questions; questions 2, 4, 8 are How questions. Note the correct order: question word + auxiliary + subject + verb.}

\tcblower

\textbf{Step 2 -- The 150-word radio report.}\par
\emph{``Good afternoon, dear listeners of Radio Balafon! Today, our team visited the maintenance workshop of the technical high school, where computers, phones and printers are repaired every day. The workshop is managed by Mr. Essono and his team of four technicians. Every week, thirty-five mobile phones, twelve laptops and eight printers are repaired here. Broken screens are the most frequent problem: they represent 40 \% of all repairs. We asked the manager: ``Which tools do you use?'' He answered that the multimeter is used for testing electric circuits, and the soldering iron is used with tin wire. USB flash drives are used for installing repair software. Surprisingly, Mr. Essono prefers repairing laptops because, he says, the work is cleaner and better paid. Six students are also trained here during their internships. This workshop shows that technical education really prepares young people for jobs. This was Youth and Jobs, on Radio Balafon!''}\par
(about 150 words)

\textbf{Step 3 -- Checking the criteria.}
\begin{itemize}
\item Passive voice (3+): \emph{are repaired}, \emph{is managed}, \emph{are trained} $\rightarrow$ correct: \emph{be} + past participle.
\item \emph{Used for / used with} (2): \emph{used for testing}, \emph{used with tin wire} $\rightarrow$ correct structures.
\item Preference (1): \emph{prefers repairing laptops} $\rightarrow$ correct.
\item Punctuation: capital letters, commas in lists, question mark after the quoted question, exclamation marks for the radio style $\rightarrow$ correct.
\end{itemize}

\textbf{Step 4 -- Remediation.} If a group wrote \emph{``the multimeter is used for test circuits''}, correct it together: after \emph{used for}, the verb takes \emph{-ing} $\rightarrow$ \emph{used for testing circuits}.
\end{exoresolu}

\courssub{Lesson Summary}

This activity has shown that the different skills of the lesson are not separate: in real life, a journalist uses them all at the same time. To produce the radio report, you had to \textbf{ask questions} (WH-/How questions), \textbf{listen and take notes} (listening), \textbf{describe machines and their uses} (passive voice, \emph{used for/with}), \textbf{give opinions} (preferences), \textbf{write a clear text} (punctuation, composition) and \textbf{speak in public} (oral presentation).

The peer-evaluation grid helped you to identify your own errors, especially the three classic difficulties: forgetting \emph{be} in the passive, forgetting \emph{-ing} after \emph{used for}, and wrong word order in questions. Keep this toolbox in mind: it will be useful for the Probatoire and Baccalauréat technique, and also in your future professional life, where technical information is often communicated through reports, interviews and the media.

\newpage

\coursec{Méthodes et exercices résolus}

\coursec{Media and Communication}

\courssub{Speaking: Exchanging Information through the Media}

\begin{definition}[Key Vocabulary -- Les médias]
\begin{center}
\begin{tabularx}{\linewidth}{|l|X|}\hline
\textbf{English} & \textbf{Français} \\ \hline
TV (Television) & Télévision \\
Radio & Radio \\
Print media / Newspaper & Presse écrite / Journal \\
Headline & Titre (une) \\
Broadcast / Programme & Émission / Programme \\
Journalist / Reporter & Journaliste / Reporter \\
Viewer / Listener / Reader & Telespectateur / Auditeur / Lecteur \\
News anchor / Presenter & Présentateur (TV) \\
Editor / Editor-in-chief & Rédacteur / Rédacteur en chef \\
Breaking news & Dernière minute / Info en direct \\
Live / Recorded & En direct / Enregistré \\
\hline
\end{tabularx}
\end{center}
\end{definition}

\begin{exemplebox}[Dialogue: Discussing News Sources]
\textbf{Context:} Two students, Paul and Marie, talk about how they get information.

\textbf{Paul:} \emph{How do you usually get your news, Marie?}

\textbf{Marie:} \emph{I prefer watching TV. The news anchor shows pictures and videos. It is easier to understand. What about you?}

\textbf{Paul:} \emph{I listen to the radio on my way to school. It is faster. Sometimes I read the newspaper at the weekend.}

\textbf{Marie:} \emph{Do you use social media?}

\textbf{Paul:} \emph{Yes, but I check the source. Fake news spreads fast on WhatsApp.}

\textbf{Marie:} \emph{True. Traditional media like TV and radio are more reliable.}
\end{exemplebox}

\begin{methode}[Useful Expressions for Exchanging Information]
\begin{itemize}
\item \textbf{Asking for information:} \emph{What is the headline today? / Did you hear the news on the radio? / Have you read the latest newspaper?}
\item \textbf{Giving information:} \emph{The TV reported that... / The newspaper says... / According to the radio...}
\item \textbf{Expressing preference (Speaking link):} \emph{I prefer TV to radio because... / I would rather listen to the radio than watch TV.}
\item \textbf{Agreeing/Disagreeing:} \emph{I agree. / That's right. / I don't think so. / Actually, the radio said...}
\end{itemize}
\end{methode}

\begin{exercice}[Speaking Practice]
Work in pairs. Role-play a conversation about yesterday's news.
\begin{enumerate}
\item Student A: Ask Student B how they got the news (TV, radio, print, online).
\item Student B: Answer and give one detail from the news.
\item Student A: React and give your own preference using \emph{prefer} or \emph{would rather}.
\item Switch roles.
\end{enumerate}
\end{exercice}

\courssub{Grammar Focus 1: Active and Passive Voices}

\begin{aretenir}[Active vs. Passive Voice -- Règles de base]
\begin{center}
\begin{tabularx}{\linewidth}{|l|X|X|}\hline
\textbf{Aspect} & \textbf{Active Voice (Voix active)} & \textbf{Passive Voice (Voix passive)} \\ \hline
\textbf{Focus} & The \textbf{doer} (subject) performs the action. & The \textbf{receiver} (object) becomes the subject. \\ \hline
\textbf{Structure} & Subject + Verb + Object & Subject + \textbf{be} (tense) + \textbf{Past Participle} (+ by + agent) \\ \hline
\textbf{Example (Present)} & The journalist \textbf{writes} the article. & The article \textbf{is written} (by the journalist). \\ \hline
\textbf{Example (Past)} & The technician \textbf{repaired} the radio. & The radio \textbf{was repaired} (by the technician). \\ \hline
\textbf{Example (Future)} & They \textbf{will broadcast} the match. & The match \textbf{will be broadcast}. \\ \hline
\textbf{Example (Pres. Perf.)} & She \textbf{has sent} the email. & The email \textbf{has been sent}. \\ \hline
\end{tabularx}
\end{center}
\textbf{Key Rule:} Only transitive verbs (verbs with a direct object / COD) can be made passive. \emph{Intransitive verbs (go, sleep, arrive, happen) cannot be passive.}
\end{aretenir}

\begin{methode}[Transforming Active to Passive -- Step by Step]
Follow these steps carefully:
\begin{enumerate}
\item \textbf{Identify} the \cle{verb} (and its tense) and the \cle{Direct Object (COD)} in the active sentence. \emph{Example: The mechanic repairs the engine. (Verb: repairs -- Present Simple; COD: the engine)}
\item \textbf{Move} the COD to the \textbf{Subject} position. \emph{The engine...}
\item \textbf{Conjugate} the auxiliary \textbf{be} in the \textbf{same tense} as the active verb. \emph{Present Simple $\rightarrow$ is/are. Past Simple $\rightarrow$ was/were. Future $\rightarrow$ will be. Present Perfect $\rightarrow$ has/have been.}
\item \textbf{Add} the \textbf{Past Participle} (Participe Passé) of the main verb. \emph{...is repaired.}
\item \textbf{Decide} on the agent (\emph{by} + doer). Keep it only if it adds important information. If the agent is unknown, obvious, or "people/they", \textbf{drop it}.
\end{enumerate}
\textbf{Reflex:} PASSIVE = \textbf{Subject + be (correct tense) + Past Participle (+ by agent)}.
\end{methode}

\begin{exoresolu}[Passive Voice -- Type Bac/Probatoire]
Rewrite the following sentences in the passive voice.\\
a) The technician repairs the radio every week.\\
b) Millions of viewers watched the match on TV last night.\\
c) The students will send the report by e-mail.\\
d) Someone has stolen my mobile phone.\\
e) They broadcast the news at 8 p.m. every day.
\tcblower
\textbf{Detailed Solutions.}

a) \emph{The technician repairs the radio every week.}
\begin{itemize}
\item Verb: \emph{repairs} (Present Simple), COD: \emph{the radio} (singular).
\item New subject: \emph{The radio}.
\item \emph{be} at Present Simple (singular) $\rightarrow$ \emph{is}.
\item Past Participle of \emph{repair} $\rightarrow$ \emph{repaired}.
\item Agent \emph{by the technician} kept (specific doer).
\end{itemize}
\textbf{Answer: The radio is repaired by the technician every week.}

b) \emph{Millions of viewers watched the match on TV last night.}
\begin{itemize}
\item Verb: \emph{watched} (Past Simple), COD: \emph{the match} (singular).
\item \emph{be} at Past Simple (singular) $\rightarrow$ \emph{was}.
\item Past Participle: \emph{watched}.
\item Agent \emph{millions of viewers} is vague/general $\rightarrow$ \textbf{drop it}.
\end{itemize}
\textbf{Answer: The match was watched on TV last night.}

c) \emph{The students will send the report by e-mail.}
\begin{itemize}
\item Verb: \emph{will send} (Future), COD: \emph{the report}.
\item Future Passive: \emph{will be} + Past Participle.
\item Past Participle of \emph{send} (irregular) $\rightarrow$ \emph{sent}.
\end{itemize}
\textbf{Answer: The report will be sent by e-mail (by the students).}

d) \emph{Someone has stolen my mobile phone.}
\begin{itemize}
\item Verb: \emph{has stolen} (Present Perfect), COD: \emph{my mobile phone}.
\item \emph{be} at Present Perfect $\rightarrow$ \emph{has been}.
\item Past Participle of \emph{steal} (irregular) $\rightarrow$ \emph{stolen}.
\item Agent \emph{someone} $\rightarrow$ \textbf{drop it}.
\end{itemize}
\textbf{Answer: My mobile phone has been stolen.}

e) \emph{They broadcast the news at 8 p.m. every day.}
\begin{itemize}
\item Verb: \emph{broadcast} (Present Simple -- same form), COD: \emph{the news}.
\item \emph{be} at Present Simple $\rightarrow$ \emph{is/are} (\emph{news} is singular in grammar $\rightarrow$ \emph{is}).
\item Past Participle of \emph{broadcast} $\rightarrow$ \emph{broadcast}.
\item Agent \emph{they} (impersonal) $\rightarrow$ \textbf{drop it}.
\end{itemize}
\textbf{Answer: The news is broadcast at 8 p.m. every day.}

\textbf{Conclusion:} Only the auxiliary \emph{be} changes tense. The Past Participle stays the same. \textbf{Watch out for irregular verbs:} \emph{send $\rightarrow$ sent, write $\rightarrow$ written, make $\rightarrow$ made, steal $\rightarrow$ stolen, broadcast $\rightarrow$ broadcast.}
\end{exoresolu}

\begin{exercice}[Passive Voice Practice]
Turn into passive. Omit the agent where natural.
\begin{enumerate}
\item The editor checks the articles before printing.
\item The company launched a new smartphone yesterday.
\item We use this machine for cutting metal.
\item They have installed the new software.
\item You must switch off the power before repairing.
\end{enumerate}
\end{exercice}

\courssub{Listening \& Vocabulary: New Information and Communication Technologies (ICT)}

\begin{definition}[ICT Vocabulary -- Vocabulaire des NTIC]
\begin{center}
\begin{tabularx}{\linewidth}{|l|X|}\hline
\textbf{Term} & \textbf{Meaning / Utilisation} \\ \hline
World Wide Web (www) / Internet & Réseau mondial / Internet \\
MP3 player & Lecteur MP3 (fichiers audio) \\
USB flash drive / Pen drive & Clé USB \\
Smartphone & Smartphone / Téléphone intelligent \\
WhatsApp / Viber / Telegram & Applications de messagerie instantanée \\
Social media / Social networks & Réseaux sociaux (Facebook, X, TikTok...) \\
Download (v) / Download (n) & Télécharger / Téléchargement \\
Upload (v) & Téléverser / Envoyer (fichier) \\
Browse / Surf the web & Naviguer sur le web \\
Wi-Fi / Connection & Connexion sans fil / Connexion \\
Software / App / Application & Logiciel / Application \\
Hardware & Matériel (physique) \\
Bluetooth & Bluetooth (connexion courte portée) \\
Cloud storage & Stockage en ligne (nuage) \\
\hline
\end{tabularx}
\end{center}
\end{definition}

\begin{experience}[Listening Simulation: "My Digital Day"]
\textbf{Script (Teacher reads aloud / Audio played):}
\emph{``Good morning. My name is David. I am a technical student. My day starts with my smartphone. I use it \textbf{as an alarm clock}. Then I check \textbf{WhatsApp} messages \textbf{with} my mobile data. In class, we use the \textbf{Internet} to research projects. I save my work on a \textbf{USB flash drive} and also back it up to \textbf{cloud storage}. In the evening, I listen to music on my \textbf{MP3 player} or stream it via \textbf{Bluetooth} on a speaker. I rarely watch \textbf{TV} now; I prefer \textbf{YouTube} or social media. Technology makes life easier, but we must protect our data.''}

\textbf{Task:} Listen and complete the table.
\begin{center}
\begin{tabularx}{\linewidth}{|l|X|}\hline
\textbf{Device / Tool} & \textbf{Function (What for?)} \\ \hline
Smartphone & Alarm clock, messaging (WhatsApp), mobile data \\ \hline
Internet / Wi-Fi & Research, browsing \\ \hline
USB Flash drive / Cloud & Storing / Backing up files \\ \hline
MP3 Player / Bluetooth & Listening to music \\ \hline
\end{tabularx}
\end{center}
\end{experience}

\begin{methode}[Expressing Purpose: Used for / Used to / Used in / Used with]
To describe the function of a gadget or machine, use these structures:
\begin{enumerate}
\item \textbf{Used for + Verb-ing} (General purpose). \emph{A USB drive is used for \textbf{storing} files.}
\item \textbf{Used to + Base Verb} (Infinitive purpose). \emph{A USB drive is used \textbf{to store} files.}
\item \textbf{Used in + Verb-ing / Noun} (Context/Field). \emph{Computers are used in \textbf{teaching} / \textbf{education}.}
\item \textbf{Used with + Noun} (Accompaniment/Requirement). \emph{WhatsApp is used \textbf{with} an Internet connection.}
\end{enumerate}
\textbf{Golden Rules:}
\begin{itemize}
\item After preposition \textbf{for / in} $\rightarrow$ \textbf{Verb + -ing} (Gerund). \emph{for storing, in teaching.}
\item After \textbf{to} (infinitive marker in \emph{used to}) $\rightarrow$ \textbf{Base Verb} (no \emph{-ing}, no \emph{to} added). \emph{used to store.}
\item After \textbf{with} $\rightarrow$ \textbf{Noun / Noun Phrase}. \emph{with a computer, with Wi-Fi.}
\item \textbf{Distinction:} \emph{Be used to + V-ing} (Habit/Accustomed) $\neq$ \emph{Used to + Base Verb} (Purpose). \emph{I am used to waking up early.} vs \emph{This tool is used to cut metal.}
\end{itemize}
\end{methode}

\begin{exoresolu}[Used for / to / in / with -- Type Bac]
Complete with the correct form of the verb in brackets or the right structure.
\begin{enumerate}
\item A smartphone is used for \dots (send) emails and \dots (browse) the web.
\item This software is used \dots (design) 3D models.
\item Robots are used in \dots (assemble) cars in modern factories.
\item A modem is used \dots (connect) a network to the Internet.
\item This machine must be used \dots (with) protective gloves.
\item My father \dots (use) to read newspapers, now he uses a tablet.
\end{enumerate}
\tcblower
\textbf{Solutions.}
\begin{enumerate}
\item \textbf{sending} / \textbf{browsing} (\emph{used for + V-ing}).
\item \textbf{to design} (\emph{used to + Base Verb}).
\item \textbf{assembling} (\emph{used in + V-ing} -- context).
\item \textbf{to connect} (\emph{used to + Base Verb}).
\item \textbf{with} (\emph{used with + Noun} -- \emph{protective gloves} is a noun phrase).
\item \textbf{used to read} (\emph{used to + Base Verb} = past habit, contrast with \emph{now}).
\end{enumerate}
\textbf{Note:} Check the word \textbf{before} the gap: \emph{for/in} $\rightarrow$ \emph{-ing}; \emph{to} $\rightarrow$ base verb; \emph{with} $\rightarrow$ noun.
\end{exoresolu}

\begin{exercice}[ICT Vocabulary \& Grammar Practice]
\begin{enumerate}
\item Match the device to its function using \emph{used for/used to/used with}.
\begin{center}
\begin{tabular}{|l|l|}\hline
\textbf{Device} & \textbf{Function} \\ \hline
1. MP3 Player & a. making video calls \\ \hline
2. Webcam & b. storing large files \\ \hline
3. External Hard Drive & c. playing audio files \\ \hline
4. Router & d. connecting multiple devices to Wi-Fi \\ \hline
\end{tabular}
\end{center}
\item Write 4 sentences describing the items above using the passive structures learned.
\end{enumerate}
\end{exercice}

\courssub{Reading: Visiting a Workshop}

\begin{definition}[Workshop Vocabulary -- Atelier / Maintenance]
\begin{center}
\begin{tabularx}{\linewidth}{|l|X|}\hline
\textbf{English} & \textbf{Français} \\ \hline
Workshop / Garage & Atelier / Garage \\
Mechanic / Technician & Mécanicien / Technicien \\
Lathe & Tour (machine) \\
Drill / Drilling machine & Perceuse / Perceuse à colonne \\
Grinding machine & Meuleuse \\
Welding machine / Welder & Poste à souder \\
Spanner / Wrench & Clé (plate, à molette, à pipe) \\
Screwdriver & Tournevis \\
Pliers & Pinces \\
Hammer & Marteau \\
Safety goggles / Glasses & Lunettes de protection \\
Gloves & Gants \\
Overall / Boiler suit & Combinaison de travail \\
Maintenance / Servicing & Maintenance / Entretien \\
Worn out / Damaged & Usé / Endommagé \\
To oil / To lubricate & Huiler / Lubrifier \\
To tighten / To loosen & Serrer / Desserer \\
To replace / To repair & Remplacer / Réparer \\
Safety regulations & Règles de sécurité \\
\hline
\end{tabularx}
\end{center}
\end{definition}

\begin{exemplebox}[Reading Text: A Visit to ENEAS Workshop]
\textbf{Last Friday,} the students of Première Technique \textbf{visited} the ENEAS mechanical workshop in Douala. The visit \textbf{was organised} by the English teacher to show how machines \textbf{are handled} and \textbf{maintained}.

\textbf{When we arrived,} the chief mechanic \textbf{welcomed} us. He \textbf{showed} us the main machines: lathes, drilling machines, and a welding station. He \textbf{explained} that safety \textbf{is} the priority. \emph{``Before any repair, the machine \textbf{must be switched off} and \textbf{unplugged},''} he \textbf{said}. \emph{``Tools \textbf{are kept} in the toolbox. Worn parts \textbf{are replaced} immediately.''}

\textbf{We learnt} the maintenance routine: every morning, the machines \textbf{are cleaned} and \textbf{oiled}. The mechanic \textbf{checked} the oil level and \textbf{tightened} loose bolts. We \textbf{observed} a technician \textbf{replacing} a damaged drill bit. He \textbf{wore} safety goggles and gloves.

\textbf{It was} a very instructive visit. We \textbf{understood} that precision and safety \textbf{are essential} in a workshop.
\end{exemplebox}

\begin{exercice}[Reading Comprehension]
Answer the questions based on the text.
\begin{enumerate}
\item When and where did the visit take place?
\item Who organised the visit and why?
\item List three machines mentioned in the text.
\item What must be done before any repair? (Use passive voice from text).
\item What is the daily maintenance routine?
\item What safety equipment did the technician wear?
\item Find in the text one sentence in the \textbf{Active Voice} (Past Simple) and one in the \textbf{Passive Voice} (Present Simple).
\end{enumerate}
\end{exercice}

\courssub{Grammar Focus 2: Expressing Preferences, Punctuation, WH- and How Questions}

\begin{methode}[Expressing Preferences: Prefer / Would Rather / Like Better]
To compare two options (media, machines, activities):
\begin{enumerate}
\item \textbf{Prefer + Gerund (-ing) + to + Gerund (-ing)} \\
\emph{I prefer \textbf{watching} TV \textbf{to} \textbf{listening} to the radio.} \\
\emph{I prefer \textbf{reading} newspapers \textbf{to} \textbf{browsing} the web.}
\item \textbf{Prefer + Noun + to + Noun} \\
\emph{I prefer \textbf{the radio} \textbf{to} \textbf{television}.}
\item \textbf{Would rather + Base Verb + than + Base Verb} \\
\emph{I \textbf{would rather} \textbf{watch} TV \textbf{than} \textbf{listen} to the radio.} \\
\emph{He'd rather \textbf{use} a lathe \textbf{than} \textbf{operate} a drill.}
\item \textbf{Like ... better than ...} \\
\emph{I \textbf{like} newspapers \textbf{better than} magazines.}
\end{enumerate}
\textbf{Traps:} \emph{Prefer} takes \textbf{to} (never \emph{than}). \emph{Would rather} takes \textbf{than} (never \emph{to}). After \emph{would rather}, verb = \textbf{Base Form} (no \emph{to}, no \emph{-ing}).
\end{methode}

\begin{exoresolu}[Expressing Preferences -- Type Bac]
\textbf{Question:} \emph{``Which do you prefer, using a smartphone or a computer for studying?''} Write two sentences (one with \emph{prefer}, one with \emph{would rather}) giving a reason.
\tcblower
\textbf{Model Answers:}
\begin{enumerate}
\item \textbf{I prefer using a computer to using a smartphone for studying because the screen is larger and typing is easier.} (\emph{Prefer + V-ing + to + V-ing})
\item \textbf{I would rather use a computer than use a smartphone because I can open multiple windows at the same time.} (\emph{Would rather + Base V + than + Base V})
\end{enumerate}
\end{exoresolu}

\begin{methode}[Punctuation Essentials -- Ponctuation Essentielle]
Correct punctuation makes writing clear. Check these rules:
\begin{enumerate}
\item \textbf{Full Stop (.)} : End of a statement. \emph{The workshop is open.}
\item \textbf{Question Mark (?)} : End of a question. \emph{When does it close?}
\item \textbf{Exclamation Mark (!)} : Strong feeling / Order. \emph{Stop the machine!}
\item \textbf{Comma (,)} :
\begin{itemize}
\item Lists: \emph{We saw lathes, drills, and grinders.}
\item After introductory words: \emph{First, switch off the power.}
\item Before connectors (\emph{but, and, so}) in long sentences.
\end{itemize}
\item \textbf{Apostrophe (')} : Contractions (\emph{don't, it's, can't}) or Possession (\emph{the mechanic's tools, the students' visit}).
\item \textbf{Capital Letters} : Start of sentence, \textbf{I}, Proper Nouns (Names, Places, Days, Months, Brands). \emph{On Monday, we visited SOPECAM in Yaoundé.}
\item \textbf{Colon (:)} : Introduces a list or explanation. \emph{We need three tools: a hammer, a wrench, and a screwdriver.}
\item \textbf{Quotation Marks (`` '')} : Direct Speech. \emph{The guide said: ``Wear your goggles.''}
\end{enumerate}
\end{methode}

\begin{exoresolu}[Punctuation -- Type Bac]
Rewrite with correct punctuation and capitals.
\emph{last monday our group visited the sopecam printing workshop in yaounde we saw offset printers computers and cutting machines the manager said please follow the safety rules}
\tcblower
\textbf{Corrected Version:}
\textbf{Last Monday, our group visited the SOPECAM Printing Workshop in Yaoundé. We saw offset printers, computers, and cutting machines. The manager said: ``Please, follow the safety rules.''}
\begin{itemize}
\item Capitals: \emph{Last, Monday, SOPECAM, Printing, Workshop, Yaoundé, We, The, Please}.
\item Commas: After \emph{Monday}, in list (\emph{printers, computers, and cutting machines}).
\item Colon + Quotes for direct speech. Full stops at end.
\end{itemize}
\end{exoresolu}

\begin{methode}[Forming WH- and How Questions]
To ask for specific information:
\begin{enumerate}
\item \textbf{Choose the Question Word:}
\begin{center}
\begin{tabular}{|l|l|}\hline
\textbf{Word} & \textbf{Information sought} \\ \hline
Who & Person (Subject/Object) \\
What & Thing / Action \\
Where & Place \\
When & Time \\
Why & Reason \\
Which & Choice \\
Whose & Possession \\
How & Manner / Quality \\
How many / much & Quantity \\
How often & Frequency \\
How long & Duration \\
\hline
\end{tabular}
\end{center}
\item \textbf{Apply Order:} \textbf{Question Word + Auxiliary + Subject + Main Verb?}
\begin{itemize}
\item Ordinary verbs: \emph{do / does / did} + Base Verb. \emph{Where \textbf{does} the mechanic \textbf{work}?}
\item Verb \emph{be} / Modals: Inversion. \emph{Where \textbf{is} the workshop? / How \textbf{can} I \textbf{repair} it?}
\item \textbf{Subject Questions (Who/What as subject):} \textbf{NO Auxiliary}. \emph{\textbf{Who} \textbf{repairs} the machines?} (Not \emph{Who does repair...})
\end{itemize}
\end{enumerate}
\end{methode}

\begin{exoresolu}[WH- / How Questions -- Type Bac]
Write the question for the underlined answer.
\begin{enumerate}
\item The workshop is \emph{in Bafoussam}.
\item The mechanic repairs \emph{ten machines} a day.
\item They visited the factory \emph{last Friday}.
\item She cleans the engine \emph{twice a week}.
\item \emph{Mr. Ngono} teaches Technical Drawing.
\item The repair costs \emph{5000 FCFA}.
\end{enumerate}
\tcblower
\textbf{Solutions.}
\begin{enumerate}
\item \textbf{Where is the workshop?} (Place $\rightarrow$ Where; Verb \emph{be} $\rightarrow$ inversion).
\item \textbf{How many machines does the mechanic repair a day?} (Quantity countable $\rightarrow$ How many; Aux \emph{does} + Base Verb \emph{repair}).
\item \textbf{When did they visit the factory?} (Time $\rightarrow$ When; Past $\rightarrow$ Aux \emph{did} + Base Verb \emph{visit}).
\item \textbf{How often does she clean the engine?} (Frequency $\rightarrow$ How often; Present $\rightarrow$ Aux \emph{does} + Base Verb \emph{clean}).
\item \textbf{Who teaches Technical Drawing?} (Subject Question $\rightarrow$ Who + Verb \emph{teaches} -- NO auxiliary).
\item \textbf{How much does the repair cost?} (Quantity uncountable/price $\rightarrow$ How much; Aux \emph{does} + Base Verb \emph{cost}).
\end{enumerate}
\end{exoresolu}

\begin{exercice}[Grammar Consolidation: Preferences, Punctuation, Questions]
\begin{enumerate}
\item Write 2 sentences comparing \emph{Radio} and \emph{TV} using \emph{prefer} and \emph{would rather}.
\item Punctuate: \emph{next friday we will go to the nkogbisson workshop the teacher said bring your notebooks}
\item Write questions for these answers:
\begin{itemize}
\item \emph{The lathe is used for shaping metal.} (Ask about function)
\item \emph{The apprentices wear blue overalls.} (Ask about clothing)
\item \emph{The training lasts six months.} (Ask about duration)
\end{itemize}
\end{enumerate}
\end{exercice}

\courssub{Writing \& Integration: Reporting a Workshop Visit}

\begin{methode}[Writing a Workshop Visit Report (120--150 words)]
Structure your composition in 3 parts:
\begin{enumerate}
\item \textbf{Introduction (When, Where, Who, Why)} \\
\emph{Last [Day/Month], the students of [Class] visited [Workshop Name] in [City] to learn how machines are handled and maintained.}
\item \textbf{Body (2 Paragraphs): Description \& Learning (Past Simple + Passive)} \\
\textbf{Para 1 (Observation):} \emph{We saw [machines]. The technician showed us how [machine] works.} \\
\textbf{Para 2 (Maintenance/Procedure -- Use Passive):} \emph{We learnt that machines are cleaned/oiled/checked daily. Worn parts are replaced. Safety gear must be worn.}
\item \textbf{Conclusion (Opinion + Preference)} \\
\emph{I enjoyed the visit because... I prefer [handling machine X] because... I hope to visit again.}
\end{enumerate}
\textbf{Connectors:} \emph{First, Then, After that, Finally, Because, However, Moreover.}
\textbf{Checklist:} Word count (120-150), Past Simple accuracy, Passive usage, Preference structure, Punctuation/Capitals.
\end{methode}

\begin{exoresolu}[Writing -- Type Bac Model Composition]
\textbf{Topic:} \emph{``A Visit to a Mechanical Workshop''} (approx. 125 words).
\tcblower
\textbf{Model Composition:}

Last month, our class of Première Technique visited a mechanical workshop in Bafoussam. The visit was organised by our English teacher to help us learn how machines are handled and maintained in a professional setting.

When we arrived, the chief mechanic welcomed us and showed us the different machines: lathes, drilling machines, and grinding machines. He explained that the machines are switched off and unplugged before any repair. Tools are kept in a locked toolbox for safety.

We learnt a lot about maintenance procedures. The machines are cleaned and oiled every morning. Worn-out parts are replaced immediately to avoid accidents. The mechanic stressed that safety goggles and gloves must be worn by every worker at all times.

I really enjoyed this visit because it was very instructive. I prefer handling small machines like drills because they are easier to control than big lathes. I hope we will visit another industrial site soon.

\textbf{Analysis:} $\approx$ 125 words. Past Simple (\emph{visited, organised, showed, explained, learnt, enjoyed, preferred}). Passive (\emph{are switched off, are kept, are cleaned, are replaced, must be worn}). Preference (\emph{prefer handling... to...}). Connectors (\emph{when, and, because, because}). Punctuation correct.
\end{exoresolu}

\begin{exercice}[Writing Task -- Exam Practice]
Write a composition of 120--150 words on: \emph{``A Visit to a Printing Workshop''}.
Include:
\begin{itemize}
\item Date, place, organiser, objective.
\item Machines seen (offset printer, cutter, computers, etc.).
\item Maintenance/Safety rules learnt (use Passive: \emph{are cleaned, are checked, must be worn}).
\item Your opinion and preference (use \emph{prefer} or \emph{would rather}).
\end{itemize}
\end{exercice}

\courssub{Integration Activities / Evaluation / Remediation}

\begin{attention}[Top 5 Errors to Avoid in Exams -- Pièges à éviter]
\begin{enumerate}
\item \textbf{Missing Auxiliary \emph{be} in Passive.} \\
Wrong: \emph{The article written by the journalist.} \\
Right: \emph{The article \textbf{is} written by the journalist.} (No \emph{be} = No Passive).
\item \textbf{Confusing \emph{Used to} structures.} \\
Wrong: \emph{It is used to storing.} \\
Right: \emph{It is used \textbf{for} storing} OR \emph{It is used \textbf{to} store.} \\
Check: \emph{for/in} $\rightarrow$ \emph{-ing}; \emph{to} $\rightarrow$ Base Verb; \emph{with} $\rightarrow$ Noun.
\item \textbf{Using \emph{than} with \emph{Prefer}.} \\
Wrong: \emph{I prefer radio than TV.} \\
Right: \emph{I prefer radio \textbf{to} TV.} (\emph{Than} only with \emph{would rather}).
\item \textbf{Keeping Past Tense after \emph{Did} in Questions.} \\
Wrong: \emph{When did you \textbf{went}?} \\
Right: \emph{When did you \textbf{go}?} (Aux \emph{did} carries the past; main verb = Base Form).
\item \textbf{Neglecting Punctuation / Capitals in Writing.} \\
Missing full stops, lowercase proper nouns (\emph{yaounde}), missing commas in lists, messy direct speech. \textbf{Proofread 2 minutes:} Capitals? Stops? Commas? Quotes?
\end{enumerate}
\end{attention}

\begin{exercice}[Integration Test -- Mixed Practice]
\textbf{Time: 30 min.}
\begin{enumerate}
\item \textbf{Grammar (Passive/ICT):} Put verbs in brackets into correct Passive tense or \emph{Used for/to/with} form.
\begin{enumerate}
\item The news \emph{(broadcast)} every evening at 8 pm.
\item A new antenna \emph{(install)} on the roof last week.
\item WhatsApp \emph{(use)} \dots \emph{(send)} voice messages.
\item This software \emph{(use)} \dots a licence key.
\item The computers \emph{(connect)} to the server right now (Present Continuous Passive).
\end{enumerate}
\item \textbf{Vocabulary:} Match the tool to the action.
\begin{center}
\begin{tabular}{|l|l|}\hline
1. Spanner & a. Measuring \\ \hline
2. Screwdriver & b. Tightening nuts/bolts \\ \hline
3. Calliper & c. Turning screws \\ \hline
4. Multimeter & d. Testing voltage/current \\ \hline
\end{tabular}
\end{center}
\item \textbf{Questions:} Write questions for the underlined parts.
\begin{enumerate}
\item The technician used a \emph{multimeter} to check the voltage.
\item The machine stops \emph{automatically} when it overheats.
\item \emph{Mr. Kouam} is the workshop supervisor.
\end{enumerate}
\item \textbf{Writing (50 words min):} You visited a radio station. Write a short paragraph (Past Simple + Passive) describing what you saw and learnt about the studio equipment.
\end{enumerate}
\end{exercice}

\begin{corrige}[Integration Test]
\begin{enumerate}
\item a) \textbf{is broadcast} (Present Simple Passive). b) \textbf{was installed} (Past Simple Passive). c) \textbf{is used for sending} / \textbf{is used to send}. d) \textbf{is used with}. e) \textbf{are being connected}.
\item 1-b, 2-c, 3-a, 4-d.
\item a) \textbf{What did the technician use to check the voltage?} b) \textbf{How does the machine stop when it overheats?} / \textbf{When does the machine stop automatically?} c) \textbf{Who is the workshop supervisor?} (Subject question).
\item \textbf{Sample Paragraph:} Last Tuesday, we visited the CRTV radio station in Yaoundé. We saw the recording studios and the control room. The sound engineer explained that microphones are tested before every broadcast. Cables are checked regularly. Headphones must be worn by the presenters. I learnt that digital editing software is used to cut and mix sounds. I preferred the control room because it is full of sophisticated technology.
\end{enumerate}
\end{corrige}

\newpage

\coursec{Exercices}

\courssub{Je vérifie mes connaissances}

\begin{exercice}[QCM : Media and ICT vocabulary]
Circle the correct answer (A, B or C).

\begin{enumerate}
\item A person who reads the news on television is a \ldots
\begin{itemize}
\item[A.] newsreader \item[B.] repairer \item[C.] listener
\end{itemize}
\item A USB flash drive is used for \ldots data.
\begin{itemize}
\item[A.] storing \item[B.] store \item[C.] stored
\end{itemize}
\item The newspaper is an example of \ldots media.
\begin{itemize}
\item[A.] broadcast \item[B.] print \item[C.] social
\end{itemize}
\item WhatsApp and Viber are used \ldots communication.
\begin{itemize}
\item[A.] for \item[B.] to \item[C.] with
\end{itemize}
\item A machine that welds metal parts in a workshop is a \ldots machine.
\begin{itemize}
\item[A.] welding \item[B.] washing \item[C.] printing
\end{itemize}
\end{enumerate}
\end{exercice}

\begin{exercice}[True or False : Active and passive voice]
Say whether each statement is \textbf{True} or \textbf{False}, and justify your answer by rewriting the sentence correctly when it is false.

\begin{enumerate}
\item ``Technicians repair 50 phones daily'' is in the passive voice.
\item ``The news is broadcast every evening by CRTV'' is in the passive voice.
\item In the passive voice, the subject performs the action.
\item The passive voice is formed with the verb \emph{to be} + past participle.
\item ``A new antenna has been installed on the roof'' is in the active voice.
\end{enumerate}
\end{exercice}

\begin{exercice}[Definitions : Media words]
Give a short definition in English (one or two sentences) for each of the following words. You may use a dictionary.

\begin{enumerate}
\item broadcast
\item headline
\item workshop
\item gadget
\item maintenance
\end{enumerate}
\end{exercice}

\begin{exercice}[WH- and How questions : quick check]
Complete each question with the correct question word : \emph{Who, What, Where, When, Why, How, How much, How many}.

\begin{enumerate}
\item \ldots\ldots do you listen to the radio ? --- In the evening.
\item \ldots\ldots repaired your phone ? --- Mr Nkeng, the technician.
\item \ldots\ldots data does this bundle contain ? --- 2 gigabytes.
\item \ldots\ldots is the workshop located ? --- Near the Douala central market.
\item \ldots\ldots machines are there in the workshop ? --- Twelve.
\end{enumerate}
\end{exercice}

\courssub{Je m'entraîne}

\begin{exercice}[From active to passive]
Rewrite the following sentences in the passive voice. Keep the same tense.

\begin{enumerate}
\item Technicians repair 50 phones daily.
\item The journalist interviews the Minister of Communication.
\item CRTV broadcasts the news at 8 p.m.
\item The students visited the workshop last week.
\item The engineer has installed a new antenna.
\item Many Cameroonians use WhatsApp every day.
\end{enumerate}
\end{exercice}

\begin{exercice}[Used for + V-ing / used with + noun]
Match each gadget (1--6) with its function (a--f), then write a complete sentence for each pair using \textbf{used for + verb-ing} or \textbf{used with + noun}.

\begin{center}
\begin{tabularx}{\linewidth}{|l|X|}
\hline
\rowcolor{popblueL} \textbf{Gadget} & \textbf{Function} \\
\hline
1. USB flash drive & a. listen to music \\
\hline
2. MP3 player & b. send instant messages \\
\hline
3. Smartphone & c. store and transfer files \\
\hline
4. Radio set & d. browse the Internet \\
\hline
5. Modem & e. listen to news and music programmes \\
\hline
6. WhatsApp & f. make video calls \\
\hline
\end{tabularx}
\end{center}

Example : \emph{A USB flash drive is used for storing and transferring files.}
\end{exercice}

\begin{exercice}[Punctuation and capital letters]
Rewrite the following paragraph with correct punctuation and capital letters.

\emph{crtv is the national television of cameroon it broadcasts news sports and cultural programmes every day many cameroonians also listen to radio stations such as crtv radio and equinoxe fm newspapers like cameroon tribune inform readers about national events do you prefer television radio or newspapers}
\end{exercice}

\begin{exercice}[Dialogue : at the repair shop]
Complete the dialogue between a customer and a repairer with the correct WH- or How questions, using the answers given.

\textbf{Customer :} Good morning. (1) \ldots\ldots\ldots\ldots ?

\textbf{Repairer :} My name is Mr Fotso.

\textbf{Customer :} (2) \ldots\ldots\ldots\ldots ?

\textbf{Repairer :} Your phone has a battery problem.

\textbf{Customer :} (3) \ldots\ldots\ldots\ldots ?

\textbf{Repairer :} The repair will cost 8\,000 FCFA.

\textbf{Customer :} (4) \ldots\ldots\ldots\ldots ?

\textbf{Repairer :} It will take two days.

\textbf{Customer :} (5) \ldots\ldots\ldots\ldots ?

\textbf{Repairer :} You can come back on Friday at 10 a.m.
\end{exercice}

\courssub{Je me prépare au Bac}

\begin{exercice}[Type Baccalauréat : Reading comprehension]
Read the text carefully and answer the questions that follow.

\textbf{A Visit to a Maintenance Workshop in Douala}

\emph{Last Tuesday, the students of Première technique of Lycée Technique de Bonabéri visited a machine maintenance workshop in the industrial zone of Douala. The visit was organised by their English teacher, Mrs Ebanga. When they arrived at 9 a.m., the chief technician, Mr Kamga, welcomed them and divided them into four groups. Each group was guided by a technician who explained how the machines work. The students saw welding machines, drilling machines and grinding machines. They were told that the machines are checked every morning before work begins. Broken parts are replaced immediately, and oil is used for lubricating the moving parts. Safety rules are strictly respected : helmets and gloves must be worn at all times. At 11 a.m., the students asked many questions about the training needed to become a technician. Mr Kamga explained that young people are trained in technical schools and that practical experience is very important. Before leaving, the students thanked the technicians. Back at school, each student wrote a short report about the visit. Most of them said they would rather work with modern machines than do office work.}

\begin{enumerate}
\item Say whether the following statements are \textbf{True} or \textbf{False}. Justify with a quotation from the text.
\begin{enumerate}
\item The visit took place in Yaoundé.
\item The students were welcomed by the English teacher.
\item The machines are checked every morning.
\item Safety rules are not respected in the workshop.
\end{enumerate}
\item Answer these questions with complete sentences.
\begin{enumerate}
\item Who organised the visit ?
\item How many groups were the students divided into ?
\item What is oil used for in the workshop ?
\item What did the students do when they returned to school ?
\end{enumerate}
\item Find in the text words or expressions that mean the same as :
\begin{enumerate}
\item damaged
\item put on
\item liked better
\end{enumerate}
\item Grammar : rewrite the following sentences from the text in the active voice.
\begin{enumerate}
\item ``The machines are checked every morning before work begins.''
\item ``Broken parts are replaced immediately.''
\end{enumerate}
\end{enumerate}
\end{exercice}

\begin{exercice}[Type Baccalauréat : Grammar and vocabulary in context]
A local newspaper published an article about a phone repair company in Bamenda. Transform the following sentences from the active voice into the passive voice, then answer the vocabulary questions.

\begin{enumerate}
\item Transform into the passive voice (keep the tense) :
\begin{enumerate}
\item Technicians repair 50 phones daily.
\item The company employs fifteen young technicians.
\item Customers bring their broken phones every morning.
\item The manager has bought new tools for the workshop.
\item The technicians will test all the repaired phones tomorrow.
\item The company offers a three-month guarantee on every repair.
\item Last year, the technicians fixed more than 10\,000 phones.
\item The secretary is writing a report about the workshop.
\end{enumerate}
\item Vocabulary : complete each sentence with one word from the box : \emph{antenna, bundle, screen, guarantee, spare parts, signal}.
\begin{enumerate}
\item The \ldots\ldots of my phone is cracked ; I cannot read my messages.
\item This repair comes with a six-month \ldots\ldots.
\item The technician replaced the old components with new \ldots\ldots.
\item I bought an internet \ldots\ldots of 1 GB for 500 FCFA.
\item The network \ldots\ldots is very weak in my village.
\end{enumerate}
\end{enumerate}
\end{exercice}

\begin{exercice}[Type Baccalauréat : Guided composition]
You visited the CRTV studios in Yaoundé with your class last month. Write a composition of about 150 words in which you :

\begin{enumerate}
\item say when and why you went there and who welcomed you ;
\item describe two or three things you saw (studios, cameras, microphones, control room) ;
\item explain what two pieces of equipment are used for, using \emph{used for + verb-ing} or \emph{used with + noun} ;
\item say which media job (newsreader, cameraman, sound engineer, journalist\ldots) you would rather do and give two reasons.
\end{enumerate}

Your composition must contain at least two sentences in the passive voice and correct punctuation. Marks will be awarded for : respect of the task, organisation of ideas, correct grammar (passive voice, preferences with \emph{would rather}), vocabulary and punctuation.
\end{exercice}



\newpage

\coursec{Corrigés des exercices}

\coursec{Media and Communication}

\courssub{I Check My Knowledge}

\begin{corrige}[Exercise 1 --- MCQ: Media and ICT Vocabulary]
\begin{enumerate}
\item \textbf{A} --- \emph{newsreader}: a person who reads the news on television or radio. A \emph{repairer} fixes machines; a \emph{listener} only listens to the radio.
\item \textbf{A} --- \emph{storing}: after \emph{used for}, the verb takes the \emph{-ing} form (\emph{used for + verb-ing}).
\item \textbf{B} --- \emph{print}: newspapers and magazines belong to the print media; radio and television belong to the broadcast media.
\item \textbf{A} --- \emph{for}: \emph{used for + noun} expresses purpose (\emph{used for communication}).
\item \textbf{A} --- \emph{welding}: a welding machine joins metal parts; a washing machine washes clothes; a printing machine prints documents.
\end{enumerate}
\end{corrige}

\begin{corrige}[Exercise 2 --- True or False: Active and Passive Voice]
\begin{enumerate}
\item \textbf{False.} The sentence is in the \emph{active} voice: the subject \emph{Technicians} performs the action. Passive form: \emph{50 phones are repaired daily (by technicians).}
\item \textbf{True.} Structure: \emph{is} (verb \emph{to be}, present simple) + \emph{broadcast} (past participle) + agent introduced by \emph{by}.
\item \textbf{False.} In the passive voice, the subject \emph{receives} the action; it is the agent (after \emph{by}) that performs it.
\item \textbf{True.} Passive voice = \emph{to be} (conjugated at the tense of the active verb) + past participle.
\item \textbf{False.} It is in the \emph{passive} voice: \emph{has been installed} (present perfect passive). Active form: \emph{They have installed a new antenna on the roof.}
\end{enumerate}
\end{corrige}

\begin{corrige}[Exercise 3 --- Definitions: Media Words]
Suggested answers (other correct formulations are acceptable):
\begin{enumerate}
\item \textbf{broadcast}: to send out programmes by radio or television; as a noun, the programme that is sent out.
\item \textbf{headline}: the title of a newspaper article, printed in large letters at the top of the page to attract the reader's attention.
\item \textbf{workshop}: a room or building where machines are used, repaired or maintained, or where things are made or repaired by hand.
\item \textbf{gadget}: a small electronic device or tool, such as a smartphone, an MP3 player or a USB flash drive.
\item \textbf{maintenance}: the regular work done to keep machines and equipment in good working condition.
\end{enumerate}
\end{corrige}

\begin{corrige}[Exercise 4 --- WH- and How Questions: Quick Check]
\begin{enumerate}
\item \emph{When} --- the answer (``In the evening'') indicates a \textbf{time}.
\item \emph{Who} --- the answer (``Mr Nkeng, the technician'') indicates a \textbf{person}, subject of the verb.
\item \emph{How much} --- \emph{data} is an \textbf{uncountable} noun.
\item \emph{Where} --- the answer indicates a \textbf{place} (``Near the Douala central market'').
\item \emph{How many} --- \emph{machines} is a \textbf{countable} noun in the plural.
\end{enumerate}
\textbf{Watch out:} do not confuse \emph{How much} (uncountable nouns: data, money, time) with \emph{How many} (countable nouns: machines, phones, students).
\end{corrige}

\courssub{I Practice}

\begin{corrige}[Exercise 5 --- From Active to Passive]
\begin{methode}[Transformation Method]
The object of the active sentence becomes the subject; the verb \emph{to be} takes the tense of the active verb; the main verb becomes a past participle; the former subject is introduced by \emph{by} (or omitted).
\end{methode}
\begin{enumerate}
\item 50 phones \textbf{are repaired} daily (by technicians). --- present simple.
\item The Minister of Communication \textbf{is interviewed} by the journalist. --- present simple.
\item The news \textbf{is broadcast} by CRTV at 8 p.m. --- present simple; irregular past participle \emph{broadcast}.
\item The workshop \textbf{was visited} by the students last week. --- past simple (\emph{visited} $\rightarrow$ \emph{was visited}).
\item A new antenna \textbf{has been installed} by the engineer. --- present perfect (\emph{has installed} $\rightarrow$ \emph{has been installed}).
\item WhatsApp \textbf{is used} by many Cameroonians every day. --- present simple.
\end{enumerate}
\textbf{Watch out:} the verb \emph{to be} must agree with the new subject (\emph{50 phones are}, \emph{the news is}) and keep the tense of the original sentence.
\end{corrige}

\begin{corrige}[Exercise 6 --- Used for + V-ing / Used with + Noun]
Matching: 1--c; 2--a; 3--f; 4--e; 5--d; 6--b.

Possible sentences:
\begin{enumerate}
\item A USB flash drive is \textbf{used for storing} and \textbf{transferring} files.
\item An MP3 player is \textbf{used for listening to} music.
\item A smartphone is \textbf{used for making} video calls.
\item A radio set is \textbf{used for listening to} news and music programmes.
\item A modem is \textbf{used for browsing} the Internet.
\item WhatsApp is \textbf{used for sending} instant messages. / WhatsApp is \textbf{used with} a smartphone.
\end{enumerate}
\textbf{Watch out:} \emph{used for} is followed by a verb in \emph{-ing} (or a noun); \emph{used with} is followed only by a noun (the tool or device combined with another). Never write \emph{used for to store}.
\end{corrige}

\begin{corrige}[Exercise 7 --- Punctuation and Capital Letters]
Corrected paragraph:

\emph{CRTV is the national television of Cameroon. It broadcasts news, sports and cultural programmes every day. Many Cameroonians also listen to radio stations such as CRTV Radio and Equinoxe FM. Newspapers like Cameroon Tribune inform readers about national events. Do you prefer television, radio or newspapers?}

Corrections made:
\begin{itemize}
\item Capital letters for proper nouns and acronyms: \emph{CRTV, Cameroon, CRTV Radio, Equinoxe FM, Cameroon Tribune}, and for the first word of each sentence.
\item Full stops at the end of the five statements.
\item Commas in the enumerations (\emph{news, sports and cultural programmes}; \emph{television, radio or newspapers}).
\item A question mark at the end of the direct question.
\end{itemize}
\end{corrige}

\begin{corrige}[Exercise 8 --- Dialogue: At the Repair Shop]
Suggested questions (other grammatically correct formulations are acceptable):
\begin{enumerate}
\item \emph{What is your name, please?} --- the answer gives a name.
\item \emph{What is the problem with my phone?} / \emph{What is wrong with my phone?} --- the answer describes the problem.
\item \emph{How much will the repair cost?} / \emph{How much is the repair?} --- the answer gives a price (8\,000 FCFA).
\item \emph{How long will it take?} --- the answer gives a duration (``two days'').
\item \emph{When can I come back (to collect it)?} --- the answer gives a date and time (``on Friday at 10 a.m.'').
\end{enumerate}
\textbf{Watch out:} respect the word order of interrogative sentences: question word + auxiliary + subject + main verb (\emph{How much will the repair cost?}, not \emph{How much the repair will cost?}).
\end{corrige}

\courssub{I Prepare for the Bac}

\begin{corrige}[Exercise 9 --- Baccalauréat Type: Reading Comprehension]
\begin{enumerate}
\item True or False (quotation from the text is compulsory for justification):
\begin{enumerate}
\item \textbf{False.} ``\ldots visited a machine maintenance workshop in the industrial zone of \emph{Douala}.''
\item \textbf{False.} ``\ldots the \emph{chief technician, Mr Kamga}, welcomed them\ldots'' (Mrs Ebanga only organised the visit).
\item \textbf{True.} ``\ldots the machines are checked every morning before work begins.''
\item \textbf{False.} ``Safety rules are \emph{strictly respected}: helmets and gloves must be worn at all times.''
\end{enumerate}
\item Answers in complete sentences:
\begin{enumerate}
\item The visit was organised by their English teacher, Mrs Ebanga.
\item The students were divided into four groups.
\item Oil is used for lubricating the moving parts.
\item When they returned to school, each student wrote a short report about the visit.
\end{enumerate}
\item Synonyms: (a) \emph{broken} (= damaged); (b) \emph{worn} (= put on); (c) \emph{would rather} (= liked better).
\item Active voice:
\begin{enumerate}
\item The technicians check the machines every morning before work begins.
\item The technicians replace broken parts immediately.
\end{enumerate}
\end{enumerate}
\textbf{Indicative marking scheme (out of 20):} question 1: 4 $\times$ 2 pts (1 pt for answer, 1 pt for quotation) = 8 pts; question 2: 4 $\times$ 1.5 pt = 6 pts; question 3: 3 $\times$ 1 pt = 3 pts; question 4: 2 $\times$ 1.5 pt = 3 pts.

\textbf{Watch out:} in question 1, an answer without a quotation from the text earns only half the points; in question 4, the tense must be preserved (present simple) and the active subject must be a logical agent (\emph{the technicians}).
\end{corrige}

\begin{corrige}[Exercise 10 --- Baccalauréat Type: Grammar and Vocabulary in Context]
\begin{enumerate}
\item Passive voice:
\begin{enumerate}
\item 50 phones \textbf{are repaired} daily (by technicians). --- present simple.
\item Fifteen young technicians \textbf{are employed} by the company. --- present simple.
\item Their broken phones \textbf{are brought} by customers every morning. --- present simple; irregular past participle \emph{brought}.
\item New tools \textbf{have been bought} for the workshop by the manager. --- present perfect; irregular past participle \emph{bought}.
\item All the repaired phones \textbf{will be tested} tomorrow (by the technicians). --- future with \emph{will}: \emph{will be} + past participle.
\item A three-month guarantee \textbf{is offered} on every repair (by the company). --- present simple.
\item Last year, more than 10\,000 phones \textbf{were fixed} (by the technicians). --- past simple; plural: \emph{were}.
\item A report about the workshop \textbf{is being written} by the secretary. --- present continuous passive: \emph{is being} + past participle; irregular past participle \emph{written}.
\end{enumerate}
\item Vocabulary: (a) \emph{screen}; (b) \emph{guarantee}; (c) \emph{spare parts}; (d) \emph{bundle}; (e) \emph{signal}.
\end{enumerate}
\textbf{Indicative marking scheme (out of 20):} question 1: 8 $\times$ 2 pts = 16 pts (1 pt for correct form of \emph{to be} in the right tense and number, 1 pt for past participle); question 2: 5 $\times$ 0.8 pt = 4 pts.

\textbf{Watch out:} check agreement of \emph{to be} with the new subject (\emph{phones were}, \emph{a guarantee is}); do not forget irregular past participles (\emph{brought, bought, written}); in the future, the form is always \emph{will be} + past participle, without agreement.
\end{corrige}

\begin{corrige}[Exercise 11 --- Baccalauréat Type: Guided Composition]
Model answer (about 150 words):

\emph{Last month, our class visited the CRTV studios in Yaoundé. The visit was organised by our English teacher so that we could learn how television programmes are made. When we arrived, we were welcomed by a journalist who guided us through the building.}

\emph{First, we entered the news studio, where the news is broadcast every evening. We saw large cameras, bright lights and many microphones. Then we visited the control room, which is full of screens and computers. We learnt that a microphone is used for recording voices, and that the cameras are used with special lights to produce clear pictures.}

\emph{I really enjoyed this visit. If I had to choose a media job, I would rather be a sound engineer, because I love technology and I would like to work with modern machines. Besides, sound engineers are needed in every radio and television station in Cameroon.}

\textbf{Indicative marking scheme (out of 20):}
\begin{itemize}
\item Task completion (all 4 points of the prompt addressed): 5 pts.
\item Organisation of ideas (paragraphs, connectors \emph{first, then, besides}): 4 pts.
\item Grammar: at least two correct passive sentences (3 pts), preference expressed with \emph{would rather} + base verb (2 pts), other correct structures (2 pts).
\item Media and ICT vocabulary: 2 pts.
\item Punctuation and capital letters: 2 pts.
\end{itemize}

\textbf{Watch out:} \emph{would rather} is followed by the base verb without \emph{to} (\emph{I would rather be}, not \emph{I would rather to be}); the passive voice requires \emph{to be} + past participle (\emph{the news is broadcast}); pay attention to punctuation in enumerations and capital letters for proper nouns (\emph{CRTV, Yaoundé, Cameroon}).
\end{corrige}

\newpage

\coursec{Fiche bilan}

\begin{popfigure}
\begin{center}
\begin{center}\tcbox[colback=popdark,colframe=popdark,coltext=white,arc=9pt,boxsep=4pt,left=6pt,right=6pt]{\bfseries\small Media and Communication}\end{center}
\vspace{-2pt}
\begin{tcbraster}[raster columns=3,raster equal height=rows,raster column skip=6pt,raster row skip=6pt,enhanced,blankest]
\begin{tcolorbox}[enhanced,colback=popblue,colframe=popblue,coltext=white,boxrule=0.9pt,arc=3pt,left=4pt,right=4pt,top=3pt,bottom=3pt,boxsep=1pt,halign=left]\bfseries\scriptsize Speaking: TV, radio, print media\end{tcolorbox}
\begin{tcolorbox}[enhanced,colback=poporange,colframe=poporange,coltext=white,boxrule=0.9pt,arc=3pt,left=4pt,right=4pt,top=3pt,bottom=3pt,boxsep=1pt,halign=left]\bfseries\scriptsize Active vs Passive voice\end{tcolorbox}
\begin{tcolorbox}[enhanced,colback=popgreen,colframe=popgreen,coltext=white,boxrule=0.9pt,arc=3pt,left=4pt,right=4pt,top=3pt,bottom=3pt,boxsep=1pt,halign=left]\bfseries\scriptsize ICT: www, MP3, USB, WhatsApp\end{tcolorbox}
\begin{tcolorbox}[enhanced,colback=poppurple,colframe=poppurple,coltext=white,boxrule=0.9pt,arc=3pt,left=4pt,right=4pt,top=3pt,bottom=3pt,boxsep=1pt,halign=left]\bfseries\scriptsize used for + V-ing / used with + noun\end{tcolorbox}
\begin{tcolorbox}[enhanced,colback=poppink,colframe=poppink,coltext=white,boxrule=0.9pt,arc=3pt,left=4pt,right=4pt,top=3pt,bottom=3pt,boxsep=1pt,halign=left]\bfseries\scriptsize Reading: visiting a workshop\end{tcolorbox}
\begin{tcolorbox}[enhanced,colback=popgold,colframe=popgold,coltext=white,boxrule=0.9pt,arc=3pt,left=4pt,right=4pt,top=3pt,bottom=3pt,boxsep=1pt,halign=left]\bfseries\scriptsize Punctuation, WH- \& How questions\end{tcolorbox}
\begin{tcolorbox}[enhanced,colback=popmuted,colframe=popmuted,coltext=white,boxrule=0.9pt,arc=3pt,left=4pt,right=4pt,top=3pt,bottom=3pt,boxsep=1pt,halign=left]\bfseries\scriptsize Writing: report a workshop visit\end{tcolorbox}
\end{tcbraster}

\end{center}
\legende{Carte mentale de la leçon : les sept compétences clés autour des médias et de la communication.}
\end{popfigure}

\begin{bilan}
\textbf{1. Key vocabulary (à connaître par cœur)}
\begin{itemize}
\item \cle{Media}: TV, radio, newspaper, magazine, headline, broadcast, presenter, newsreader, advert(isement).
\item \cle{ICT}: website, to browse, to download, to upload, MP3 player, USB flash drive, password, to chat, network.
\item \cle{Workshop}: machine, tool, gadget, engine, to repair, to maintain, maintenance, spare part, safety rules.
\end{itemize}

\textbf{2. Grammar must-knows}
\begin{center}
\begin{tabularx}{\linewidth}{|l|X|X|}\hline
\rowcolor{popblueL} \textbf{Point} & \textbf{Rule} & \textbf{Example} \\ \hline
Active voice & Subject + verb + object & The journalist \textbf{writes} the article. \\ \hline
Passive voice & Object + \textbf{be} + past participle (+ by...) & The article \textbf{is written} by the journalist. \\ \hline
Used for & \textbf{used for + V-ing} (purpose) & A USB key is used \textbf{for storing} files. \\ \hline
Used with & \textbf{used with + noun} & This software is used \textbf{with} a computer. \\ \hline
Preferences & prefer, would rather, like... better than & I \textbf{prefer} radio \textbf{to} TV. \\ \hline
WH-/How questions & Wh-word + auxiliary + subject + verb? & \textbf{How often} do you listen to the radio? \\ \hline
\end{tabularx}
\end{center}

\textbf{3. Punctuation reminders}
\begin{itemize}
\item Full stop (.) ends a sentence; question mark (?) after a question; comma (,) in lists and after linkers; capital letter at the start of sentences and for proper nouns; apostrophe (') for contractions and possession.
\end{itemize}

\textbf{4. Express methods}
\begin{itemize}
\item \cle{Active $\rightarrow$ Passive}: move the object to subject position, add \emph{be} in the same tense, put the verb in the past participle, add \emph{by} + agent if needed.
\item \cle{Workshop report plan}: date and place $\rightarrow$ machines observed $\rightarrow$ operations explained (use the passive!) $\rightarrow$ preferences and feelings $\rightarrow$ conclusion.
\item \cle{Asking questions}: choose the right WH-word (who, what, where, when, why, how) + correct auxiliary (do/does/did/is/are).
\end{itemize}
\end{bilan}

\begin{aretenir}[Checklist avant le Bac]
\begin{itemize}
\item $\square$ I can name at least 10 words about the media and 10 about ICT.
\item $\square$ I can transform an active sentence into the passive voice in the present and past simple.
\item $\square$ I never forget \emph{be} + past participle in the passive.
\item $\square$ I can use \emph{used for + V-ing} and \emph{used with + noun} correctly.
\item $\square$ I can express preferences with \emph{prefer ... to} and \emph{would rather}.
\item $\square$ I can build WH- and How questions with the correct word order.
\item $\square$ I can punctuate a short paragraph correctly (capital letters, full stops, commas, question marks).
\item $\square$ I can understand a short listening text about new technologies.
\item $\square$ I can read a text about a workshop visit and answer comprehension questions.
\item $\square$ I can write a 120-word composition reporting a workshop visit, using the passive voice.
\item $\square$ I check my spelling of key words: \emph{broadcast, maintenance, gadget, download}.
\end{itemize}
\end{aretenir}

\findecours

\end{document}
